% !TEX root = ../../mixed-sklyanin-cubics.tex
% Current result: 2026-11-22 / v53.
\section{中心二重根の全 coherent 分解と実残余関手}
\label{mix:sec:root-order}

零係数の対合を取ると、混合代数は中心元の二乗根として記述できる。
本節ではその記述から tensor と inflation の関手を作り、
全 coherent 射影圏と残余圏まで分解を降ろす。
まず可換四様体の Kuznetsov 成分から
$\operatorname{Ku}_{\mathrm{coh}}(A_0)^{\mu_2}$ への実際の完全忠実関手と、
その補成分・貼合せの具体的な記述を得る。
次に貼合せに現れる球面制限を境界上の twist 射から比較し、
定理\ref{cmo:thm:fixed-zero-commutative-kernel}で
通常の $\operatorname{Ku}_{\mathrm{coh}}(A_0)$ と指定可換候補の
$\C$-線型 dg Morita 同値の perfect 核を構成する。
この固定値の構成は、交換子パラメータの一般点核の特殊化を用いない。

可換な根スタックの半直交分解と、分岐因子の残余圏を mutation して
同変残余圏へ移す方法は既知である
\cite[Sections~4 and~6]{CMOKuznetsovPerryCovers}。
以下では本研究の非可換不変環にも適用できるよう、
idempotent と長さ一の解消から必要な関手を直接構成する。
一般的な根構成そのものの新規性は主張しない。

\subsection{次数付き order の核と全 coherent 射影圏}

本節では右加群を用い、内部 shift は $M(a)_q=M_{q+a}$ とする。
$R$ を標数が二でない体 $k$ 上の、左右 noetherian な
非負次数付き locally finite 代数とし、$R_0=k$ とする。
$v\in Z(R)_2$ は左右の非零因子と仮定する。
中心変数 $u$ の次数を一として
\[
 A=R[u]/(u^2-v),\qquad
 \sigma|_R=1,\quad \sigma(u)=-u,\qquad
 C=A\rtimes\langle\sigma\rangle,\qquad S=R/(v)
\]
と置く。$e_\pm=(1\pm\sigma)/2$、
$\mathcal T_Q=D^b(\operatorname{qgr}Q)$ と書く。
$\rho$ は $S$ 加群を $R$ 加群と見る inflation、
$\mathcal L=-\otimes_R^{\mathbf L}S$ は導来制限とする。

\begin{proposition}[次数付き order の全 coherent 分解]
\label{cmo:prop:graded-root-order}
次の関手は次数付き有限生成加群の有界導来圏上で定義され、
有限次元 cohomology を持つ複体を保存する。
従って同じ式により $\mathcal T_Q$ 上へ降りる。
\[
 \begin{aligned}
 i(N)&=N \quad\text{($C\twoheadrightarrow S$ による inflation)},\\
 i^*(Y)&=Y\otimes_C^{\mathbf L}S,
 &i^!(Y)&=\underline{R\operatorname{Hom}}_C(S,Y),\\
 j(M)&=M\otimes_R e_+C,
 &j^*(Y)&=Ye_+,\\
 j_*(M)&=\underline{\operatorname{Hom}}_R(Ce_+,M)
       \cong M(1)\otimes_R e_-C,
 &j^L(Y)&=(Ye_-)(1).
 \end{aligned}
\]
ここで underline は内部次数も保持する Hom を表す。
随伴は
\[
 i^*\dashv i\dashv i^!,\qquad
 j^L\dashv j\dashv j^*\dashv j_*
\]
である。$i,j,j_*$ は完全忠実であり、
\begin{equation}
\label{cmo:eq:root-order-coherent-sod}
 \mathcal T_C=\langle i\mathcal T_S,j\mathcal T_R\rangle.
\end{equation}
同じ分解は次数付き有界導来圏、および通常の環の perfect 圏でも成立する。
貼合せは自然な dg 同型
\begin{equation}
\label{cmo:eq:root-order-gluing}
 R\operatorname{Hom}_{\mathcal T_C}(iN,jM)
 \simeq R\operatorname{Hom}_{\mathcal T_R}(\rho N(1),M)
 \simeq R\operatorname{Hom}_{\mathcal T_S}
             (N,\mathcal L M(1))[-1]
\end{equation}
で与えられる。
\end{proposition}

\begin{proof}
次数を忘れた $R$ 加群 $A=R\oplus uR$ 上の作用により
\[
 C\cong\begin{pmatrix}R&vR\\R&R\end{pmatrix},\qquad
 Ce_+C\cong\begin{pmatrix}R&vR\\R&vR\end{pmatrix},
 \qquad e_+Ce_+\cong R,\quad C/Ce_+C\cong S.
\]
次数付きの計算では両 off-diagonal 成分を $uR\cong R(-1)$ として扱う。
行列表示の $vR$ を次数を保たずに $R$ と同一視しない。
$v$ の正則性により積写像
$Ce_+\otimes_R e_+C\to Ce_+C$ は同型であり、
左右の $R$ 加群は有限自由なので導来 Tor はない。
従って次数付き双加群の完全列
\begin{equation}
\label{cmo:eq:root-order-bimodule-sequence}
 0\longrightarrow Ce_+\otimes_R e_+C
  \longrightarrow C\longrightarrow S\longrightarrow0
\end{equation}
を得る。

右 $C$ 加群 $S$ の解消は
\begin{equation}
\label{cmo:eq:root-order-resolution}
 0\longrightarrow e_+C(-1)
  \xrightarrow{\ (e_-ue_+)\cdot\ }e_-C
  \longrightarrow S\longrightarrow0
\end{equation}
である。二つの idempotent 成分では、写像はそれぞれ
$R(-1)\to uR$ の同型と $uR(-1)\to R$ の像 $vR$ を与える。
左加群にも対応する解消がある。
この解消を $S$ と tensor すると $S\otimes_C^{\mathbf L}S\simeq S$ となるので、
$i^*i\simeq\operatorname{id}$ である。
また $j^*j\simeq\operatorname{id}$、$j^*i=0$ は idempotent の計算である。
従って $i,j$ は完全忠実で、後者から前者への全 Hom が消える。
\eqref{cmo:eq:root-order-bimodule-sequence} を右 $C$ 複体 $Y$ と tensor した三角
\begin{equation}
\label{cmo:eq:root-order-projection-triangle}
 j j^*Y\longrightarrow Y\longrightarrow i i^*Y
\end{equation}
が生成性を与える。ここで二つの射は、それぞれ積と商写像から来る実際の随伴射である。

$e_+C$ と $Ce_+$ は対応する $R$ 側で
$R\oplus R(-1)$ である。従って $j,j^*,j_*$ は exact である。
$j_*M$ の二つの成分は $(M,M(1))$ であり、
$M(1)\otimes_R e_-C$ と同じ $u$ 作用を持つ。
この同型と符号指標による $e_+,e_-$ の交換から
$j^L(Y)=(Ye_-)(1)$ を得る。
解消\eqref{cmo:eq:root-order-resolution} とその左版により、
$i^*$ と $i^!$ の cohomological amplitude はそれぞれ
$[-1,0]$ と $[0,1]$ に含まれる。
以上の全関手は、有限個の idempotent 成分、内部 shift、
有限生成加群の核・余核だけで cohomology を計算できるため、
有限生成性と有限次元 cohomology を保存する。

ここで使う有界導来 Serre 商の同一視も確認する。
本仮定の次数付き有限生成加群では torsion は有限次元加群である。
$\operatorname{qgr}Q$ の射は、十分高い次数で切った加群の間の次数零写像で表せる。
有限個の対象・射・等式からなる diagram は、一つの共通の高次数で切れば
$\operatorname{gr}Q$ の diagram に持ち上がる。
実際、射の表示に現れる有限次元の核・余核をまず除き、
等式の差の有限次元像をさらに高次数で除けばよい。
この操作は noetherian 性により有限生成性を保つ。
有界複体、roof、および有界な homotopy はいずれも有限 diagram なので、
これを適用すると
\[
 D^b(\operatorname{gr}Q)/D^b_{\mathrm{fd}}(\operatorname{gr}Q)
 \xrightarrow{\sim}D^b(\operatorname{qgr}Q)
\]
を得る。従って上の関手・随伴射・三角は全てこの商へ降り、
\eqref{cmo:eq:root-order-coherent-sod} が全 coherent 圏で成立する。
perfect 圏への制限は、$S$ が左右の $C$ 加群として perfect であることと
$e_+C,Ce_+$ の有限自由性から従う。

最後に $j^L iN=\rho N(1)$ である。
$0\to R(-2)\xrightarrow{v}R\to S\to0$ の双対により、
inflation $\rho$ の右随伴は
\[
 \underline{R\operatorname{Hom}}_R(S,M)
 \simeq (M\otimes_R^{\mathbf L}S)(2)[-1]
\]
となる。これら二つの随伴を合成すれば
\eqref{cmo:eq:root-order-gluing} を得る。
これは内部次数を忘れた Hom の数値的一致ではなく、貼合せ双加群の自然な同定である。
\end{proof}

\begin{remark}[全 coherent 圏と perfect 核の区別]
一般の $R$ に対する上の命題は、全 coherent 圏の分解を主張するが、
各圏の smoothness は仮定していない。
一方、以下の指定混合族では $\mathcal T_C$ は smooth proper である。
$j$ は左右の随伴を持つので、$\mathcal T_R$ もその admissible な dg 部分圏として
smooth proper となる。$\mathcal T_S$ にも同じことがいえる。
この場合、以下の tensor、inflation、有限 mutation の関手は
smooth proper dg 圏間の perfect 双加群核で与えられる。
一般の環に対する片側 perfect 性だけを、双加群核の perfect 性と取り違えない。
\end{remark}

% !TEX root = ../../mixed-sklyanin-cubics.tex
\subsection{不変環の重み付き周囲代数と五本の例外線束}

以下では指定した二ブロックの $c=0$ を扱い、
$w=\Omega_1+100\Omega_2$、$A_0=B/(w)$ とする。
六つの active 変数を固定し $u$ を $-u$ に送る対合を $\sigma$ と書く。
\[
 D=B^\sigma,\qquad v=u^2,\qquad
 R=A_0^\sigma=D/(w),\qquad S=R/(v)=A_0/(u)
\]
と置く。$v$ の内部次数は二である。
不変環の議論では、次数一の生成元の個数から global dimension を推定しない。
実際、$v$ は表示上は次数二の中心生成元だが、
$\alpha_j^{(i)}\ne0$ の関係式から六つの次数一生成元の二次式としても書ける。

\begin{proposition}[不変環の正則周囲代数と残余圏]
\label{cmo:prop:invariant-weighted-algebra}
$D$ は両側 noetherian な graded Calabi--Yau 代数であり、
global dimension は七、Gorenstein parameter は八である。
その Nakayama 自己同型は恒等であり、
\begin{equation}\label{cmo:eq:invariant-hilbert}
 H_D(t)=\frac{1}{(1-t)^6(1-t^2)},\qquad
 H_R(t)=\frac{1-t^3}{(1-t)^6(1-t^2)}
\end{equation}
である。$w\in D_3$ は正則中心元で、$R$ は左右 injective dimension 六、
Gorenstein parameter 五の AS-Gorenstein 代数となる。
その canonical 双加群は $\omega_R=R(-5)$ である。

$C=A_0\rtimes\langle\sigma\rangle$、$\chi$ を非自明な指標とすると、
全 coherent 導来圏の Serre 関手は
\begin{equation}\label{cmo:eq:invariant-serre}
 \mathsf S_{D^b(\operatorname{qgr}C)}=(-4)\otimes\chi[5],
 \qquad
 \mathsf S_{D^b(\operatorname{qgr}R)}=(-5)[5]
\end{equation}
である。$D^b_{\mathrm{dg}}(\operatorname{qgr}R)$ は smooth proper で、
\begin{equation}\label{cmo:eq:invariant-residual}
 D^b_{\mathrm{dg}}(\operatorname{qgr}R)
 =\langle\mathcal K_R,\OO_R,\OO_R(1),\ldots,\OO_R(4)\rangle,
 \qquad
 \mathcal K_R\simeq\MF^{\mathbb Z}(D,w)
\end{equation}
という実際の dg 分解・同値を持つ。
有限射影行列分解の右辺は冪等完備化して読む。
内部 shift に対応する $\rho_R$ と Serre 関手は
\begin{equation}\label{cmo:eq:invariant-fractional-serre}
 \rho_R^3\simeq[2],\qquad
 \mathsf S_{\mathcal K_R}\simeq\rho_R^{-5}[5],\qquad
 \mathsf S_{\mathcal K_R}^{3}\simeq[5]
\end{equation}
を満たす。最後の式は三次三様体の Ku と同じ fractional Serre 関係だが、
可換三次三様体への同値をそれだけから結論しない。
\end{proposition}

\begin{proof}
$u$ は $B$ 上正則で、$B/(u)$ 上の対合は恒等である。
反不変元 $b$ は $u$ で割り切れ、$b=ur$ と書いたとき $r$ は不変となる。
従って、左右 $D$-加群として
\[
 B=D\oplus uD,\qquad D/(v)\cong B/(u)
     \cong S(1,-1,-1)\otimes S(1,-1,-1)
\]
である。$D$ は六つの次数一生成元と $v$ で生成され、$v$ は中心正則元である。
$B$ は $D$ 上両側自由なので、$D$ の左イデアル $I$ について
$BI=I\oplus uI$、特に $BI\cap D=I$ である。
$B$ の noetherian 性から左イデアルの上昇列は停止し、右側も同様である。
商の表示と $v$ の正則性から式\eqref{cmo:eq:invariant-hilbert}の第一式を得る。

一因子 $S(1,-1,-1)$ の関係式は cyclic potential
\[
 xyz-xzy-\frac{x^3+y^3+z^3}{3}
\]
の三つの cyclic derivative である。
既知の長さ三の Koszul 双加群分解は、この cyclic 表示によって
その双対と内部次数三・長さ三の shift を除いて一致する。
具体的には外側の二写像は生成元の左右差で、中央の写像は potential の
二階 cyclic derivative であり、転置と左右交換により互いに移る。
AS-regular 性によるこの Koszul 分解の完全性と合わせて、
一因子は Nakayama twist のない graded Calabi--Yau 三次元代数である。
二つの分解を tensor すると $D/(v)$ は graded Calabi--Yau 六次元となる。

正次数の正規拡大に関する
\cite[Theorem~3.1 and Remark~2.2]{CMOZhouShenLu}を $v$ に適用する。
$D$ は連結局所有限、$\deg v=2>0$ なので同定理の正次数条件を満たし、
$D$ は skew Calabi--Yau 七次元、parameter $6+2=8$ となる。
その Nakayama 自己同型 $\mu_D$ は $v$ を法として恒等である。
$(v)$ に次数一の元はないため、$\mu_D$ は六つの次数一生成元を固定する。
さらに \cite[Proposition~2.3]{CMOZhouShenLu}により中心元 $v$ 自身も固定する。
従って $\mu_D=1$ である。

$w$ は $B$ 上の正則中心元であり $D$ に属するから、$D$ 上も正則である。
不変部分を取る関手の完全性により $A_0^\sigma=D/(w)$ である。
中心元の二項解消と Rees の補題
\cite[Lemma~3.4]{CMOZhouShenLu}は、$R$ の左右有限 injective dimension 六、
\[
 R\underline{\operatorname{Hom}}_R(\C,R)\simeq\C(5)[-6]
\]
および canonical 双加群 $R(-8+3)=R(-5)$ を与える。
同じ中心商により Hilbert 級数の第二式を得る。
noetherian な正則中心商なので、導来 section の $\chi$ 条件と有限
cohomological dimension も満たす
\cite[Proposition~4.18 and Lemmas~6.1--6.2]{LP-BS}。

Serre 関手の指標も確認する。まず $B/(u)$ の上記 Calabi--Yau 分解と
正規拡大定理から、$B$ は parameter 七の skew Calabi--Yau 代数である。
その Nakayama 自己同型は $u$ を固定し、$u$ を法として恒等である。
連結次数付き代数の Nakayama 自己同型は次数付き自己同型群の中心に属する
\cite[Proposition~2.3]{CMOZhouShenLu}ので、$\sigma$ と可換する。
従って不変な次数一生成元に反不変な $u$ の倍を足すことはできず、
この Nakayama 自己同型も恒等である。
命題\ref{cmo:prop:equivariant-window-hom}の最高次 Koszul 項は
内部次数七の一次元空間で、$\sigma$ の指標は $-1$ である。
$w$ は不変であるから、二項中心商の双対を取ってもこの指標は変わらない。
従って $A_0$ の同変 canonical 双加群は $A_0(-4)\otimes\chi$ である。
既知の全 coherent smooth proper 性と局所双対性により、
$C$ の Serre 関手は式\eqref{cmo:eq:invariant-serre}の第一式となる。

命題\ref{cmo:prop:graded-root-order}の完全忠実関手
\[
 j(M)=M\otimes_R e_+C,\qquad j^*(N)=Ne_+
\]
とその両随伴により、$D^b(\operatorname{qgr}R)$ は
$D^b(\operatorname{qgr}C)$ の admissible 成分である。
従ってこれも smooth proper であり、その Serre 関手は
$j^*\mathsf S_Cj$ で計算できる。
指標 twist は $e_+C$ と $e_-C$ を交換し、
$e_-Ce_+=uR\cong R(-1)$ なので、得られる関手は
$M\mapsto M(-5)[5]$ である。
これにより第二式を、coherent 対象を取り落とさずに得た。

AS-Gorenstein の局所 cohomology から
\[
 \operatorname{Ext}^q_{\operatorname{qgr}R}
       (\OO_R,\OO_R(d))=
 \begin{cases}
 R_d,&q=0,\\
 R_{-d-5}^{\vee},&q=5,\\
 0,&q\ne0,5
 \end{cases}
\]
である。負の添字の次数成分は零と読む。
従って $\OO_R,\ldots,\OO_R(4)$ は例外列となる。
Burke--Stevenson の Orlov 分解
\cite[Theorem~6.4]{LP-BS}を parameter 五で適用し、
有限大域次元の $D$ と正則中心元 $w$ の余核比較を用いると、
式\eqref{cmo:eq:invariant-residual}を得る。
この余核比較は、行列分解の二つの写像の合成が指定元 $w$ であることから
内部 shift の周期 $(3)\simeq[2]$ を保つ。

最後に、この同じ window 比較で一段の内部 shift を $\rho_R$ と書く。
Serre 関手 $(-5)[5]$ を残余に射影するとき、五本の例外線束を横切る
五回の window 変更が内部 shift $(-5)$ に対応する。
従って残余 Serre 関手は $\rho_R^{-5}[5]$ となる。
これは各一段の変更を評価三角による mutation で実現した式であり、
残余の外部 twist $(-5)$ への安定性を仮定していない。
三乗して $(3)\simeq[2]$ を使えば $[-10+15]=[5]$ を得る。
\end{proof}


\subsection{八つの例外対象を除く実際の mutation}

以下は指定された $A_0$、$R=A_0^\sigma$、$S=A_0/(u)$ に戻る。
$\mathcal O_R(r)=\pi R(r)$、$\mathcal O_S(r)=\pi S(r)$ と書き、
\[
 \mathcal K_S=\langle\mathcal O_S,\mathcal O_S(1),\mathcal O_S(2)\rangle^\perp,
 \qquad
 \mathcal K_R=\langle\mathcal O_R,\ldots,\mathcal O_R(4)\rangle^\perp
\]
と置く。これらは全 coherent 圏の残余部分圏である。
$C$ 側の二種類の線束を
\[
 E_r^+=\pi e_+C(r)=j\mathcal O_R(r),\qquad
 E_r^-=\pi e_-C(r)
\]
と書く。$C$ 加群と $\sigma$ 同変 $A_0$ 加群の対応の下で、
これらは $\mathcal O_{A_0}(r)$ の自明・符号 linearization に対応する。
有限群の位数が可逆なので、同変化は torsion 商と有界導来圏に両立する
\cite[Section~8]{MEDElagin}。従って
\[
 \mathcal K_C:=
 \langle E_0^+,E_0^-,E_1^+,E_1^-,E_2^+,E_2^-,E_3^+,E_3^-\rangle^\perp
 \simeq\operatorname{Ku}_{\mathrm{coh}}(A_0)^{\mu_2}.
\]

\begin{theorem}[同変残余圏の実二成分分解]
\label{cmo:thm:root-order-residual}
関手
\begin{equation}
\label{cmo:eq:root-order-boundary-kernel}
 \begin{split}
 \Phi(F)&=L_{E_0^-}\bigl(i(F(1))\bigr)\\
 &=\operatorname{Cone}\bigl(
   R\operatorname{Hom}_{\mathcal T_S}(\mathcal O_S,F(1))
      \otimes_k E_0^-
    \longrightarrow i(F(1))\bigr),\qquad
 \Psi(M)=jM
 \end{split}
\end{equation}
はそれぞれ $\mathcal K_S,\mathcal K_R$ から $\mathcal K_C$ への
perfect 核を持つ完全忠実 dg 関手であり、
\begin{equation}
\label{cmo:eq:root-order-residual-sod}
 \operatorname{Ku}_{\mathrm{coh}}(A_0)^{\mu_2}
 \simeq\mathcal K_C
 =\langle\Phi\mathcal K_S,\Psi\mathcal K_R\rangle.
\end{equation}
さらに $\mathcal L=-\otimes_R^{\mathbf L}S$ は
$\mathcal K_R\to\mathcal K_S$ を与え、貼合せは
\begin{equation}
\label{cmo:eq:root-order-residual-gluing}
 R\operatorname{Hom}_{\mathcal K_C}(\Phi F,\Psi M)
 \simeq R\operatorname{Hom}_{\mathcal K_S}(F,\mathcal L M)[-1].
\end{equation}
特に $\Phi^R\Psi\simeq\mathcal L[-1]$ である。
\end{theorem}

\begin{proof}
命題\ref{cmo:prop:invariant-weighted-algebra}により $R$ の parameter は五であり、
$S$ の parameter は三である。従って
\[
 \mathcal T_S=
 \langle\mathcal K_S(1),\mathcal O_S(1),\mathcal O_S(2),\mathcal O_S(3)\rangle,
 \quad
 \mathcal T_R=\langle\mathcal K_R,\mathcal O_R,\ldots,\mathcal O_R(4)\rangle.
\]
まず $1\le r\le3$ と $M\in\mathcal K_R$ に対して
\[
 R\operatorname{Hom}_{\mathcal T_C}(i\mathcal O_S(r),jM)
 \simeq R\operatorname{Hom}_{\mathcal T_R}(\rho\mathcal O_S(r+1),M)=0
\]
である。実際、$\rho\mathcal O_S(r+1)$ の解消は
$[\mathcal O_R(r-1)\xrightarrow{v}\mathcal O_R(r+1)]$ であり、
二つの index はともに $0,\ldots,4$ に入る。
逆方向の Hom は\eqref{cmo:eq:root-order-coherent-sod}により零である。
従って三つの $i\mathcal O_S(r)$ は $j\mathcal K_R$ と交換でき、
\[
 \mathcal T_C=
 \langle i\mathcal K_S(1),j\mathcal K_R,
   i\mathcal O_S(1),i\mathcal O_S(2),i\mathcal O_S(3),
   j\mathcal O_R,\ldots,j\mathcal O_R(4)\rangle.
\]

$\mathcal T_C$ の Serre 関手を $\mathbb S_C$ とする。
左右の随伴の一意性により
$\mathbb S_Cj\simeq j_*\mathbb S_R$ である。
$\mathbb S_R=(-5)[5]$ と命題\ref{cmo:prop:graded-root-order}の $j_*$ の式から
\[
 \mathbb S_C(j\mathcal O_R(4))
 \simeq j_*\mathcal O_R(-1)[5]
 \simeq E_0^-[5].
\]
最後の例外対象を Serre 関手により先頭へ回すと
\begin{equation}
\label{cmo:eq:root-order-rotated-sod}
 \mathcal T_C=
 \langle E_0^-,i\mathcal K_S(1),j\mathcal K_R,
    i\mathcal O_S(1),i\mathcal O_S(2),i\mathcal O_S(3),
    E_0^+,E_1^+,E_2^+,E_3^+\rangle.
\end{equation}
cohomological shift は生成する部分圏に影響しない。

次に
\[
 R\operatorname{Hom}_{\mathcal T_C}(E_0^-,jM)
 \simeq R\operatorname{Hom}_{\mathcal T_R}(\mathcal O_R(1),M)=0
 \quad(M\in\mathcal K_R)
\]
なので、先頭の $E_0^-$ を二つの残余成分の右へ移す left mutation は
$j\mathcal K_R$ には作用せず、$i\mathcal K_S(1)$ には
\eqref{cmo:eq:root-order-boundary-kernel}として作用する。
ここで
$R\operatorname{Hom}_{\mathcal T_C}(iN,E_0^-)=0$ は、
$E_0^-=j_*\mathcal O_R(-1)$ と $j^*i=0$ から全ての $N$ について成立する。
従ってこの mutation は $i\mathcal K_S(1)$ 上で完全忠実である。
明示的には、evaluation triangle と
$R\operatorname{Hom}(iN,E_0^-)=R\operatorname{Hom}(E_0^-,L_{E_0^-}iN)=0$
を二度用いれば、mutation 前後の全 Hom が同型になる。

この操作で右に残る八つの例外対象は
\begin{equation}
\label{cmo:eq:root-order-eight-mutated-exceptionals}
 E_0^-,i\mathcal O_S(1),i\mathcal O_S(2),i\mathcal O_S(3),
 E_0^+,E_1^+,E_2^+,E_3^+
\end{equation}
である。解消\eqref{cmo:eq:root-order-resolution}の shift は三角
\[
 E_{r-1}^+\xrightarrow{u}E_r^-
 \longrightarrow i\mathcal O_S(r),\qquad r=1,2,3
\]
を与えるので、\eqref{cmo:eq:root-order-eight-mutated-exceptionals}の生成部分圏は
元の八つの $E_r^\pm$、$0\le r\le3$ のものに等しい。
これで全八条件への直交性、生成性、
\eqref{cmo:eq:root-order-residual-sod} が証明された。
半直交分解の数値的な相殺は用いていない。

$M\in\mathcal K_R$ と $0\le r\le2$ に対し、
$\rho\mathcal L M$ は
$[M(-2)\xrightarrow{v}M]$ で表される。
この制限は構造対象からの導来 Hom に両立するため、
$R\operatorname{Hom}_{\mathcal T_S}(\mathcal O_S(r),\mathcal L M)$ は
\[
 R\operatorname{Hom}_{\mathcal T_R}
     (\mathcal O_R(r),[M(-2)\xrightarrow{v}M])
\]
と同型である。両項は $\mathcal O_R(r+2),\mathcal O_R(r)$ から $M$ への Hom なので零である。
従って $\mathcal L M\in\mathcal K_S$ となる。
最後に\eqref{cmo:eq:root-order-boundary-kernel}の triangle に
$R\operatorname{Hom}(-,jM)$ を適用する。
$R\operatorname{Hom}(E_0^-,jM)=0$ なので correction term が消え、
\eqref{cmo:eq:root-order-gluing}の二つの内部 shift $(1)$ が相殺して
\eqref{cmo:eq:root-order-residual-gluing} を得る。
\end{proof}

\begin{corollary}[可換四様体成分からの実関手]
\label{cmo:cor:commutative-boundary-embedding}
$Z=V(G)\subset\mathbb P^5$ を命題\ref{cmo:prop:zero-eckardt-resolution}の
分離型三次四様体とする。既存の分離比較と
\eqref{cmo:eq:root-order-boundary-kernel}の合成は、
perfect 核を持つ完全忠実関手
\[
 \operatorname{Ku}(Z)\longrightarrow
       \operatorname{Ku}_{\mathrm{coh}}(A_0)^{\mu_2}
\]
を与える。
\end{corollary}
\begin{proof}
$S=(S(1,-1,-1)\otimes S(1,-1,-1))/(\Omega_1+100\Omega_2)$ である。
二つの楕円因子の偏極点曲線は
$x^3+y^3+z^3+xyz=0$ であるから、
主稿の実因子に対する分離比較\paperref{main}{pps:thm:elliptic-existence}により
$\mathcal K_S\simeq\operatorname{Ku}(Z)$ を得る。
これを定理\ref{cmo:thm:root-order-residual}の $\Phi$ と合成すればよい。
これは同変圏への埋込みであり、通常の五様体残余圏への完全忠実埋込みを主張しない。
\end{proof}

\subsection{残余制限関手の spherical 性}

貼合せに現れた導来制限は、任意の exact functor ではない。
次の命題では cotwist の shift 規約も固定する。

\begin{proposition}[導来制限の cotwist と随伴射]
\label{cmo:prop:invariant-restriction-spherical}
定理\ref{cmo:thm:root-order-residual}の
$F=\mathcal L|_{\mathcal K_R}:\mathcal K_R\to\mathcal K_S$ は
perfect 核を持つ spherical dg functor である。
$\varrho_R=L_{\mathcal O_R}\circ(1)|_{\mathcal K_R}$ を rotation とすると、
\begin{equation}
\label{cmo:eq:invariant-restriction-cotwist}
 \mathsf C_F:=\operatorname{Cone}(\operatorname{id}\longrightarrow F^RF)
 \simeq\varrho_R^{-2}[1]
 \simeq\mathbb S_{\mathcal K_R}[-2].
\end{equation}
特に $\operatorname{Cone}(FF^R\to\operatorname{id}_{\mathcal K_S})$ は自己同値である。
\end{proposition}

\begin{proof}
包含を $a:\mathcal K_R\hookrightarrow\mathcal T_R$、
その随伴を $p_L\dashv a\dashv p_R$ と書く。
$\rho:\mathcal T_S\to\mathcal T_R$ は先に定めた inflation である。
全圏の制限 $\mathcal L$ の右随伴は $\rho$、左随伴は $\rho(2)[-1]$ なので、
\[
 F^R=p_R\rho|_{\mathcal K_S},\qquad
 F^L=p_L\rho(2)[-1]|_{\mathcal K_S}.
\]
中心元 $v$ の解消を tensor した随伴三角
\[
 M(-2)\xrightarrow{v}M\longrightarrow\rho\mathcal L M
\]
へ $p_R$ を適用すれば、
$\mathsf C_F(M)\simeq p_R(M(-2))[1]$ を得る。

inverse rotation の実 mutation 式は
\[
 \varrho_R^{-1}=R_{\mathcal O_R(-1)}\circ(-1)|_{\mathcal K_R}.
\]
二回適用すると、自然射
$a\varrho_R^{-2}M\to M(-2)$ の cone は
$\langle\mathcal O_R(-2),\mathcal O_R(-1)\rangle$ に属する。
この二対象は $\mathcal K_R^\perp$ に属する。
実際、$\mathbb S_{\mathcal T_R}=(-5)[5]$ による Serre 双対で、
それぞれ $\mathcal O_R(3),\mathcal O_R(4)$ から $\mathcal K_R$ への消滅に帰着する。
従って $p_R$ はその cone を殺し、
$p_R(M(-2))\simeq\varrho_R^{-2}M$ となる。
これで\eqref{cmo:eq:invariant-restriction-cotwist}の第一の同型と
$\mathsf C_F$ の可逆性を得た。
第二の同型は\eqref{cmo:eq:invariant-fractional-serre}の
$\varrho_R^3\simeq[2]$ と
$\mathbb S_{\mathcal K_R}=\varrho_R^{-5}[5]$ による。

spherical 性には、随伴を同定する標準射そのものを確認する。
$N\in\mathcal K_S$ に対して $Y=\rho N(2)[-1]$ と置く。
全圏の中心 Cartier 制限について、長さ一の解消とその双対が与える標準射
\[
 \rho N\longrightarrow ((-2)[1])\bigl(\rho N(2)[-1]\bigr)
\]
は同型である。これは解消の生成元とその双対による評価であり、
同一視後は $\rho N$ の恒等写像となる。
制限関手 $F$ の単位を合成すると、標準射
$F^R N\to\mathsf C_F F^L N$ は
\begin{equation}
\label{cmo:eq:invariant-restriction-adjoint-map}
 p_R\bigl(Y(-2)[1]\bigr)
 \longrightarrow p_R\bigl((a p_LY)(-2)[1]\bigr)
\end{equation}
となる。右の射は $p_L\dashv a$ の単位 $Y\to a p_LY$ から誘導される。
この記述は中心元の解消から作った単位・余単位を合成したものであり、
二つの随伴の抽象的な Serre 同型を代入したものではない。

$2\le r\le4$ について
\[
 R\operatorname{Hom}_{\mathcal T_R}(\mathcal O_R(r),Y)
 \simeq R\operatorname{Hom}_{\mathcal T_S}(\mathcal O_S(r-2),N)[-1]=0.
\]
従って $Y\to a p_LY$ を五つの線束の left mutation で計算すると、
非零の correction term は $\mathcal O_R(0),\mathcal O_R(1)$ だけからなる。
その cone に $(-2)[1]$ を施すと
$\langle\mathcal O_R(-2),\mathcal O_R(-1)\rangle$ に入り、
$p_R$ により零になる。
従って\eqref{cmo:eq:invariant-restriction-adjoint-map}は同型である。

以上により可逆な cotwist と標準随伴射の同型を確認したので、
\cite[Theorem~5.1, conditions~(2) and~(4)]{CMOAnnoLogvinenkoSpherical}
を適用できる。同文献の cotwist は本稿の $\mathsf C_F[-1]$ であり、
標準随伴射はそこで $F^R\to(\mathsf C_F[-1])F^L[1]$ と書かれる。
全ての cone は既に構成した dg 双加群核の cone である。
従って $F$ は spherical であり、記した twist は自己同値となる。
\end{proof}

\subsection{指標作用と核を比較するための残りの入力}

指標 $\chi$ による作用は、次数付き代数自己同型
\[
 \varphi(r)=r,\qquad\varphi(u)=u,\qquad
 \varphi(\sigma)=-\sigma
\]
で右作用を twist する関手である。
その核は ${}_CC_\varphi$ であり、$\varphi^2=1$ により
$\chi^2\simeq\operatorname{id}$ の coherence も指定される。
この関手は $e_+$ と $e_-$ を交換し、内部 shift と可換で、
八つの $E_r^\pm$ の集合を保存する。従って $\mathcal K_C$ 上に作用する。
この作用を残余分解の二成分ごとの恒等作用とみなしてはいけない。
実際、$j_-(M)=M\otimes_R e_-C$ と書くと
\begin{equation}
\label{cmo:eq:root-order-character-triangles}
 \begin{gathered}
 \chi jM=j_-M=j_*M(-1),\\
 jM(-1)\xrightarrow{u}j_-M\longrightarrow i\mathcal L M,\\
 iN(-1)\longrightarrow j\rho N\longrightarrow\chi iN
 \end{gathered}
\end{equation}
は実際の tensor 三角である。
最後の三角は、$vN=0$ のときの $j\rho N$ の二つの成分
$(N,N(-1))$ を調べれば得られる。
従って特に
\[
 \chi\Phi(F)=L_{E_0^+}\bigl(\chi i(F(1))\bigr),\qquad
 \chi\Psi(M)=j_-M
\]
であり、\eqref{cmo:eq:root-order-character-triangles}は
これらの対象を元の order の核から再構成する。

忘却関手 $\operatorname{For}$ については
\[
 R\operatorname{Hom}_{\mathcal T_{A_0}}
    (\operatorname{For}X,\operatorname{For}Y)
 \simeq R\operatorname{Hom}_{\mathcal T_C}(X,Y)
       \oplus R\operatorname{Hom}_{\mathcal T_C}(X,\chi Y)
\]
となる。従って上の完全忠実性から、忘却後の完全忠実性は従わない。
通常の $\operatorname{Ku}$ を回復するには、$\chi$ とその二乗の coherence を保った
比較を作り、二重同変化\cite[Theorem~4.2 and Section~8]{MEDElagin}を用いる。
本例では、この coherence は同変残余圏の dg 同値が得られた後に補正できる。

\begin{proposition}[指標作用の Serre 関手による回復]
\label{cmo:prop:root-order-intrinsic-character}
$\mathcal K_C$ の Serre 関手を $\mathbb S_{\mathcal K_C}$ とすると、
\begin{equation}
\label{cmo:eq:root-order-intrinsic-character}
 \chi\simeq\mathbb S_{\mathcal K_C}^{3}[-7].
\end{equation}
可換候補の $\operatorname{Ku}(X_0)^{\mu_2}$ にも同じ式が成立する。
従って両者の dg 同値は、指標作用を自然同型の範囲で必ず保つ。
\end{proposition}

\begin{proof}
同変な Orlov/window 比較で、内部 shift に対応する残余 rotation を
\[
 \varrho_C=L_{\langle E_0^+,E_0^-\rangle}\circ(1)|_{\mathcal K_C}
\]
と書く。指定した三次中心元 $w$ は対合で不変であるため、
同変行列分解の周期は $\varrho_C^3\simeq[2]$ であり、
この同型に追加の符号指標は入らない。
指標 twist は二本の例外対象を交換するので $\varrho_C$ と可換する。
式\eqref{cmo:eq:invariant-serre}の同変 canonical 双加群を
四段の window 変更に通すと
\[
 \mathbb S_{\mathcal K_C}\simeq\varrho_C^{-4}\chi[5]
\]
を得る。これは $\mathcal K_R$ の Serre 式と同じ評価三角による計算で、
各段に二つの指標の例外対象がある点だけが異なる。
従って
\[
 \mathbb S_{\mathcal K_C}^{3}
 \simeq\varrho_C^{-12}\chi^3[15]
 \simeq\chi[7].
\]
可換側でも周囲の七変数に対する最高次指標は $-1$、
三次方程式の指標は $+1$ なので、全く同じ計算が適用される。
同値は Serre 双対性を保ち、Serre 関手を自然同型の範囲で同定するため、
最後の主張が従う。
\end{proof}

\begin{lemma}[二乗の coherence の平方根補正]
\label{cmo:lem:root-order-coherence-square-root}
$\mathcal U,\mathcal V$ を $\C$ 上の smooth proper dg 圏とし、
それぞれに二乗の coherence を指定した対合 $\tau,\tau'$ が作用するとする。
両作用には strict な dg Morita 模型が取れると仮定する。
perfect 核 $P$ による dg 同値 $H:\mathcal U\to\mathcal V$ と、
核の導来圏における同型から来る自然同型 $H\tau\simeq\tau'H$ が与えられれば、
この自然同型を補正し、perfect 核を実際に同変化することで、
$\mu_2$ 同変圏の dg Morita 同値を構成できる。
\end{lemma}

\begin{proof}
核 $P$ の次数零自己準同型代数を $Z$ とする。
$P$ は可逆な Morita 核なので
$Z\cong HH^0(\mathcal U)$ であり、これは有限次元可換 $\C$ 代数である。
与えた自然同型を二回合成し、指定された二乗の同型と比較した不整合は
ある $d\in Z^\times$ で表される。
より具体的には、核の圏上の対合
$Q\mapsto\tau'Q\tau^{-1}$ と $P$ 上の同型を用いて
$Z$ 上に自己同型 $\gamma$ を定めると、
二回の合成は $d$ による共役となり、三回の合成の結合則は
$\gamma(d)=d$ を与える。
$Z$ が可換なので $\gamma^2=1$ である。

有限次元可換 $\C$ 代数の単元 $d$ には、$d$ の多項式で表される
平方根 $b=d^{-1/2}$ がある。
実際、$d$ の相異なる非零固有値で局所分解し、各固有値の平方根を選ぶ。
各成分の冪零部分では binomial 展開は有限和となり、
中国剰余定理で一つの多項式にまとめられる。
従って $\gamma(b)=b$ であり、$b\gamma(b)d=b^2d=1$ となる。
元の自然同型を $b$ で補正すれば二乗の不整合が消える。

これは自然変換の間の等式として、
\cite[Definition~3.1]{MEDElagin}の群作用の coherence を満たす。
従って equivariant object の構造写像へ補正後の同型を合成することで
同変圏の関手が実際に得られる。逆同値にも対応する構造を移せるため同値である。
元の関手は exact であり、同変圏の distinguished triangle は忘却後に検出される
\cite[Definition~6.1 and Corollary~6.10]{MEDElagin}ので、この同値も exact である。

次に核を dg 圏上で同変化する。
strict な作用を持つ生成子の dg 自己準同型代数を
$\Gamma_{\mathcal U},\Gamma_{\mathcal V}$ とし、
$\mathscr E=\Gamma_{\mathcal U}^{\mathrm{op}}\otimes\Gamma_{\mathcal V}$ と置く。
$G=\mu_2$ の対角作用は $\mathscr E$ 上 strict であり、
補正後の核は $D(\mathscr E)$ 内で
$\theta:P\xrightarrow{\sim}gP$、$g(\theta)\theta=1$ を持つ。
strict な同変 dg 加群については
\[
 R\operatorname{Hom}_{\mathscr E\rtimes G}(Q,Q')
 \simeq R\operatorname{Hom}_{\mathscr E}(Q,Q')^G
\]
である。実際、$\mathscr E\rtimes G$ は $\mathscr E$ 上自由なので
cofibrant 解消の制限も cofibrant であり、Hom 複体の等式をその解消上で計算できる。
不変部分を取る関手は平均作用素により exact なので、
\[
 D(\mathscr E\rtimes G)\longrightarrow D(\mathscr E)^G
\]
は完全忠実である。

$\operatorname{Ind}(P)=P\oplus gP$ には成分を交換する実際の strict 同変構造がある。
これは誘導加群であるから $\mathscr E\rtimes G$ 上 perfect である。
その導来自己準同型
\begin{equation}
\label{cmo:eq:root-order-equivariant-kernel-projector}
 e=\frac12
 \begin{pmatrix}1&\theta^{-1}\\\theta&1\end{pmatrix}
 \in\operatorname{End}_{D(\mathscr E)}(P\oplus gP)
\end{equation}
は $G$ 不変な idempotent で、その像は弱同変対象 $(P,\theta)$ である。
上の完全忠実性によって $e$ を
$\operatorname{Perf}(\mathscr E\rtimes G)$ の idempotent へ持ち上げ、
この冪等完備な圏で分裂させれば、実際の strict 同変 perfect 核 $\widetilde P$ を得る。
この構成は有限群同変 enhancement の
\cite[Theorems~8.7 and~8.9]{MEDElagin}を、核の dg 加群圏で直接実現したものである。
個々の高次 homotopy を選び足す必要はない。

$\widetilde P$ を対角部分群 $G\subset G\times G$ から誘導すれば、
$\Gamma_{\mathcal U}\rtimes G$ と $\Gamma_{\mathcal V}\rtimes G$ の間の
perfect 双加群核を得る。
これは同変加群を $P$ と tensor し、対角の同変構造を入れる関手そのものである。
右随伴は双対核に誘導される同変構造で与えられる。
その単位・余単位を忘却すると元の Morita 同値の単位・余単位となり、いずれも同型である。
同変 dg 加群の忘却は conservative なので、これらは同変圏でも同型となる。
従って構成した核は dg Morita 同値を与える。
\end{proof}

本例の $\chi$ は元の skew group 代数の自己同型 $\varphi$ による twist であり、
torsion の dg 商と八つの例外対象の直交部分圏を保存する。
従って補題の strict 模型の仮定を満たす。
具体的には生成対象を $G_0\oplus\chi G_0$ に置き換え、成分交換を用いれば、
自己準同型 dg 代数上に二乗が恒等の strict な作用を得られる。

\begin{proposition}[貼合せ核からの構成的比較]
\label{cmo:prop:root-order-comparison-criterion}
可換候補の不変環と境界環を
\[
 R^{\mathrm{cl}}=
  \C[x_1,\ldots,x_6,v]/(G-vL),\qquad
 \deg x_i=1,\quad\deg v=2,\qquad
 S^{\mathrm{cl}}=R^{\mathrm{cl}}/(v)
\]
とする。上と同じ残余部分圏と関手を上付き $\mathrm{cl}$ で表す。
perfect 核を持つ dg 同値
\[
 F_S:\mathcal K_S\xrightarrow{\sim}\mathcal K_{S^{\mathrm{cl}}},
 \qquad
 F_R:\mathcal K_R\xrightarrow{\sim}\mathcal K_{R^{\mathrm{cl}}}
\]
と、各成分が quasi-isomorphism である dg 自然変換
\begin{equation}
\label{cmo:eq:root-order-restriction-intertwiner}
 F_S\mathcal L\xrightarrow{\sim}\mathcal L^{\mathrm{cl}}F_R
\end{equation}
が与えられれば、これらの核を貼り合わせることで
\[
 \widehat F:\operatorname{Ku}_{\mathrm{coh}}(A_0)^{\mu_2}
       \xrightarrow{\sim}\operatorname{Ku}(X_0)^{\mu_2}
\]
を実際に構成できる。
指標への両立を別の入力として課さずに、
$\widehat F$ の核を平方根補正・同変化して二重同変化すれば、
通常の残余圏の perfect 核を持つ dg Morita 同値
$\operatorname{Ku}_{\mathrm{coh}}(A_0)\simeq\operatorname{Ku}(X_0)$ を得る。
\end{proposition}

\begin{proof}
二成分分解\eqref{cmo:eq:root-order-residual-sod}は、
対角成分が $\mathcal K_S,\mathcal K_R$、一方向の Hom が零、
他方向の Hom が\eqref{cmo:eq:root-order-residual-gluing}である dg 貼合せである。
その二つの対角上でそれぞれ
\[
 \Phi F\longmapsto\Phi^{\mathrm{cl}}F_SF,
 \qquad
 \Psi M\longmapsto\Psi^{\mathrm{cl}}F_RM
\]
と定め、他方向の Hom を
\eqref{cmo:eq:root-order-restriction-intertwiner}によって写す。
dg 自然性は対角との合成を保つので、これは dg 貼合せの関手を与える。
有限 twisted complex と冪等完備化へ延長すると $\widehat F$ が得られる。
対角の全 Hom と貼合せ双加群上で quasi-isomorphism であり、両成分を生成するため同値である。
同じ構成は、二つの対角核と貼合せ写像を持つ三角行列の双加群核として表せる。
従って単に同値の存在を要求する判定ではなく、入力核から比較核を構成している。
単位と余単位は両成分の生成対象上で同型であり、cone と直和因子を取って全圏で同型になる。
命題\ref{cmo:prop:root-order-intrinsic-character}により
$\widehat F\chi\simeq\chi^{\mathrm{cl}}\widehat F$ がある。
補題\ref{cmo:lem:root-order-coherence-square-root}で二乗の coherence を合わせ、
指定された $\mu_2$ 作用について二重同変化
\cite[Theorem~4.2]{MEDElagin}を適用する。
この最後の操作も dg Morita 同値であることは、strict な dg 代数模型 $\Gamma$ に対する
\[
 (\Gamma\rtimes\mu_2)\rtimes\mu_2^\vee
 \cong\operatorname{Mat}_2(\Gamma)
\]
から直接確認できる。実際、$a\in\Gamma$ を
$\operatorname{diag}(a,\sigma(a))$、最初の群の生成元を成分交換行列、
双対群の生成元を $\operatorname{diag}(1,-1)$ に送ると、
積と微分を保つ上記の同型になる。
従って通常の残余圏の間にも実際の perfect 比較核が得られる。
\end{proof}

この命題の境界核 $F_S$ は既存の分離比較から得られる。
後続の球面制限の比較では、その rotation と twist 射への両立を確認し、
全 coherence を持つ comonad の同型から $F_R$ と
\eqref{cmo:eq:root-order-restriction-intertwiner}を構成する。
指標への両立とその二乗の coherence は、この命題の補正・同変化で処理する。
従って定理\ref{cmo:thm:fixed-zero-commutative-kernel}では、
二つの対角成分の実核と貼合せ双加群の比較を経て通常の残余圏の同値を得る。
定理\ref{cmo:thm:root-order-residual}は、その構成を
不変環の残余圏と一つの具体的な導来制限関手へ還元する段階である。

% !TEX root = ../../mixed-sklyanin-cubics.tex
\subsection{可換側の実 Kn\"orrer 核と単純な超平面切断の障害}

不変環の残余 $\mathcal K_R$ を可換三次三様体と比較することは、
元の同変五様体の比較核を構成するための具体的な部分問題となる。
まず可換側では、この三様体成分を実際の核で取り出せる。

\begin{proposition}[可換不変環の次数付き Kn\"orrer 核]
\label{cmo:prop:commutative-invariant-knorrer}
\S\ref{cmo:subsec:eckardt-root-model}の $G,L,Y$ に対し、
\[
 D_{\mathrm{cl}}=\C[x_1,\ldots,x_6,v],\qquad
 \deg x_j=1,\quad\deg v=2,\qquad
 R_{\mathrm{cl}}=D_{\mathrm{cl}}/(G-vL)
\]
と置く。$Y=V(G,L)\subset\mathbb P^5$ が滑らかであるとする。
五本の例外線束の残余を $\mathcal K_{R_{\mathrm{cl}}}$ とすると、
\begin{equation}\label{cmo:eq:commutative-invariant-knorrer}
 \Ku(Y)\simeq\MF^{\mathbb Z}(\C[y_1,\ldots,y_5],g)
 \xrightarrow{\ \sim\ }
 \MF^{\mathbb Z}(D_{\mathrm{cl}},G-vL)
 \simeq\mathcal K_{R_{\mathrm{cl}}}
\end{equation}
という実 dg Morita 同値がある。
ここで $\ell=L$ を第六の線型座標に取り
\[
 G(y,\ell)=g(y)+\ell Q(y,\ell),\qquad t=v-Q(y,\ell)
\]
と置くと、中央の関手は $g-\ell t$ に対する rank 一の二変数核
\begin{equation}\label{cmo:eq:commutative-knorrer-kernel}
 \mathsf Q=\left(
  \C[\ell,t](-1)\xrightarrow{\ell}\C[\ell,t]
       \xrightarrow{-t}\C[\ell,t](2)\right),
 \qquad (\deg\ell,\deg t)=(1,2)
\end{equation}
との外積で与えられる。内部次数と指定された三次周期を保つ。
\end{proposition}

\begin{proof}
$\ell=L$ は非零線型形式なので座標に含められる。
多項式の除法により $Q$ は重み二で、$v\mapsto t=v-Q$ は
次数付き多項式環の自己同型である。この変更後の potential は
$G-vL=g-\ell t$ となる。
式\eqref{cmo:eq:commutative-knorrer-kernel}の二つの合成は
$-\ell t$ の乗法であり、両写像は表示した内部 shift に関して次数零である。
外積の微分を $d_E\otimes1+(-1)^{\deg} \otimes d_{\mathsf Q}$ と取れば、
その二乗は $g-\ell t$ となる。

同変 Kn\"orrer 周期性 \cite[Theorem~1.2]{CMOHiranoKnorrer}を用いる。
底を $\mathbb A^5_y\times\mathbb A^1_\ell$、群を $\mathbb G_m$、
potential の指標を重み三、正則切断を $\ell$ とし、
その双対変数 $t$ に重み二を与える。
切断の零集合は $\mathbb A^5_y$、制限 potential は $g$ である。
$Y$ の滑らかさから $g$ は非零であり、
$\C[z]\to\C[y]$、$z\mapsto g$ は単射で、その上の加群は torsion-free、
従って平坦である。このため同定理の制限 potential の平坦性条件も満たす。
同定理の関手は零切断への pullback と pushforward で与えられ、
その二項 Koszul 解消を取ると上の外積核になる。
ambient は正則なので coherent factorization は有限局所自由模型で計算でき、
この比較は有限射影行列分解の冪等完備後の全圏に及ぶ。

最後に $\C[y]/(g)$ の Gorenstein parameter は $5-3=2$、
$R_{\mathrm{cl}}$ の parameter は $(6+2)-3=5$ である。
それぞれに Orlov 分解を適用して両端の同値を得る。
重み付き環 $R_{\mathrm{cl}}$ の qgr は対応する重み付き射影スタックの
coherent 圏であり、粗空間へ置き換えて stacky な点を捨てていない。
外積と座標変更は potential を厳密に保つので、三次周期も保存される。
\end{proof}

非可換側で同じ線型形式を零とする操作は、上の比較に代わらない。
この点は点集合だけでなく、二側商全体について証明できる。

\begin{proposition}[指定混合不変環の線型二側商]
\label{cmo:prop:invariant-linear-quotient-finite}
$\alpha^{(1)}=(1,4,9)$、$\alpha^{(2)}=(16,25,49)$、$q_1=1$、$q_2=100$ とし、
\[
 L=\sum_{i=1}^2q_i\bigl(
   \alpha^{(i)}_1x_i+\alpha^{(i)}_2y_i+\alpha^{(i)}_3z_i\bigr)\in D_1
\]
と置く。$(L)$ は $L$ が生成する二側イデアルを表す。
このとき $D/(L)$ と $R/(L)$ は有限次元であり、いずれも $v^{33}$ で消される。
特に
\begin{equation}\label{cmo:eq:invariant-linear-empty-qgr}
 \operatorname{qgr}(R/(L))=0.
\end{equation}
従って $R$ 内の $L$ は normal 元ではなく、
この二側商を可換三様体 $Y$ の代わりに使うことはできない。
\end{proposition}

\begin{proof}
$Q=D/(L)$ と置く。異なるブロックの生成元は可換である。
$[L,x_i]=[L,y_i]=[L,z_i]=0$ と元の関係式から
\[
 \alpha^{(i)}_3y_i^2=\alpha^{(i)}_2z_i^2,\quad
 \alpha^{(i)}_1z_i^2=\alpha^{(i)}_3x_i^2,\quad
 \alpha^{(i)}_2x_i^2=\alpha^{(i)}_1y_i^2
\]
を得る。全ての $\alpha^{(i)}_j$ が非零なので
\[
 T_i=\frac{x_i^2}{\alpha^{(i)}_1}
     =\frac{y_i^2}{\alpha^{(i)}_2}
     =\frac{z_i^2}{\alpha^{(i)}_3},\qquad c_i=T_i-v
\]
を定められる。$T_i$ は各生成元との可換性をその平方表示から確認できるため中心であり、
\begin{equation}\label{cmo:eq:linear-quotient-relations}
 [y_i,z_i]=\alpha^{(i)}_1c_i,\quad
 [z_i,x_i]=\alpha^{(i)}_2c_i,\quad
 [x_i,y_i]=\alpha^{(i)}_3c_i
\end{equation}
である。例えば $[x_i^2,y_i]=2\alpha^{(i)}_3c_ix_i=0$ を使い、
他の二変数にも同じ計算をすると
\[
 c_ix_i=c_iy_i=c_iz_i=0,\qquad
 c_iT_i=0,\qquad T_i^2=vT_i
\]
となる。さらに $c_i[x_i,y_i]=\alpha^{(i)}_3c_i^2=0$ より $c_i^2=0$、
$c_i^2=(T_i-v)c_i=-vc_i$ より $vc_i=0$ である。

式\eqref{cmo:eq:linear-quotient-relations}の右辺は中心元である。
従って語を並べ替え、平方を $\alpha^{(i)}_jT_i$ で置き換え、
$T_i^2=vT_i$ を使うことにより、$Q$ は $\C[v]$-加群として
\[
 T_1^{a_1}T_2^{a_2}
 x_1^{b_1}y_1^{b_2}z_1^{b_3}
 x_2^{b_4}y_2^{b_5}z_2^{b_6},
 \qquad a_i,b_j\in\{0,1\}
\]
という $256$ 個の元で生成される。

$I=(c_1,c_2)\subset Q$ とすると $vI=0$ である。
$Q/I$ では全生成元が可換で、それぞれの平方は $\alpha^{(i)}_jv$ となる。
まず $L$ の関係を課す前の可換環
\[
 H=\C[x_1,y_1,z_1,x_2,y_2,z_2,v]/
 (x_i^2-\alpha^{(i)}_1v,\ y_i^2-\alpha^{(i)}_2v,\ z_i^2-\alpha^{(i)}_3v)_{i=1,2}
\]
を取る。$H$ は $\C[v]$ 上自由で階数 $64$ であり、$L$ の乗法の行列式は
\begin{equation}\label{cmo:eq:linear-quotient-norm}
 \det(m_L)=v^{32}
 \prod_{\epsilon\in\{\pm1\}^6}
 \bigl(\epsilon_1+8\epsilon_2+27\epsilon_3
       +100(64\epsilon_4+125\epsilon_5+343\epsilon_6)\bigr).
\end{equation}
実際、$v$ を可逆化して $s^2=v$ を添加すれば、$H$ は符号選択による
$64$ 個の一次元成分へ分かれ、各固有値は括弧内の数の $s$ 倍となる。
忠実平坦なこの拡大で得た行列式の式は元の $\C[v]$ 上にも成立する。
括弧内の最初の三項の絶対値は高々 $36$、
残りの和の絶対値は $100(343-125-64)=15400$ 以上である。
従って全ての因子は非零である。
乗法行列の adjugate 恒等式、または Cayley--Hamilton により、
$H/(L)$ はこの非零定数倍の $v^{32}$ で消される。
$Q/I$ は $H/(L)$ の商なので $v^{32}Q\subset I$、従って $v^{33}Q=0$ である。
$256$ 個の生成元による有限性と合わせて $Q$ は有限次元となる。
$R/(L)$ はその商なので同じ結論を満たす。

最後に、もし $L$ が $R$ の normal 元なら、
$R(-1)\xrightarrow{\cdot L}R\to R/(L)\to0$ から
\[
 \dim(R/(L))_n\geq\dim R_n-\dim R_{n-1}
\]
となる。式\eqref{cmo:eq:invariant-hilbert}により右辺は十分大きい $n$ で正であり、
有限次元性に矛盾する。
\end{proof}

\paragraph{有限曲率模型と球面制限による比較.}
可換側の命題\ref{cmo:prop:commutative-invariant-knorrer}は実際の比較核を与える。
非可換側の命題\ref{cmo:prop:invariant-linear-quotient-finite}が排除するのは、
同じ $L$ の通常の二側商による置換である。
$-vL$ の非退化な重み $(2,1)$ の対を curved Koszul 双対や
$A_\infty$ 模型で導来的に消去する可能性までは排除しない。
実際、定理\ref{cmo:thm:finite-curved-invariant-model}では、
元の交換関係と曲率 $w$ を保つ生成行列分解と、
停止性を証明した総階数八十への有限縮約を構成する。
この行列模型を可換三次三様体の模型へ直接移す非可換 Kn\"orrer 双加群は、
別途、全微分・高次積・随伴と unit、counit を含めて構成すべき課題である。

一方、圏の比較そのものは
定理\ref{cmo:thm:fixed-zero-commutative-kernel}の球面制限の方法により得られる。
そこでは境界の実同値と twist 射の比較から
$\mathcal K_R\simeq\Ku(Y)$ および二成分を結ぶ制限 intertwiner を同時に構成し、
通常の五様体の perfect 比較核へ貼り合わせる。
従って残る行列模型の直接比較は、この核をさらに具体的な係数で記述するための課題であり、
指定された固定値での $\C$-線型 dg Morita 同値の未証明箇所ではない。

% !TEX root = ../../mixed-sklyanin-cubics.tex
\subsection{固定係数での有限行列分解と曲率を保つ縮約}
\label{cmo:subsec:curved-invariant-model}

本節では $h=1$ において $\mathcal K_R$ の生成対象を具体化する。
前節の通常の二側商 $R/(L)$ を用いず、中心三次元 $w$ の高次 homotopy を
有限自由解消に入れる。この構成には形式パラメータの特殊化を必要としない。
得られるのは $D$ 上の有限階数の行列分解とその実際の縮約である。
その全自己準同型 dg 代数が有限次元ベクトル空間である、あるいは
既に可換三次三様体の代数と同定された、という主張ではない。

\paragraph{有限の非斉次 Koszul 解消.}
$S_0=S(1,-1,-1)$ とし、$E_0=S_0^!$ を二次双対とする。
$E_0$ は次数一の $a,b,c$ により生成され、
\[
 ab+ba=bc+cb=ca+ac=0,\qquad
 a^2+bc=b^2+ca=c^2+ab=0
\]
を関係とする。次数別次元は $(1,3,3,1)$ であり、
$1,a,b,c,a^2,b^2,c^2,T=abc$ が基底になる。
$a^3=b^3=c^3=-T$、$E_0^{\ge4}=0$ で、二次部分は中心である。

$F=E_0\mathbin{\widehat\otimes}E_0$ を符号付き tensor 積とし、
\begin{equation}\label{cmo:eq:curved-ambient-dual}
 E=F\mathbin{\widehat\otimes}\Lambda(\epsilon),\qquad
 d_E|_F=0,\qquad
 d_E\epsilon=q=-\sum_{i=1}^2\sum_{j=1}^3
        \alpha_j^{(i)}\xi_{ij}^2
\end{equation}
と置く。ここで $\epsilon$ の cohomological 次数は一、内部重みは二で、
$\xi_{ij}$ の二つの次数はいずれも一である。
$q$ は中心なので $d_E^2=0$ となる。
これは、七生成元の補助語長を一とした $D$ の quadratic--linear 双対である。
関係式 $r+\varphi(r)=0$ に対して $d_E=\varphi^*$ とする規約は
\cite[\S\S2.1--2.5]{CMOPositselskiQL}に従う。

計算の行列は左 $D$-加群で表示する。$K_i=D\otimes(E^i)^\vee$ とし、
$\eta\in(E^i)^\vee$ に対して
\begin{equation}\label{cmo:eq:curved-koszul-differential}
 d_K(1\otimes\eta)=
 \sum_{i,j}x_{ij}\otimes L_{\xi_{ij}}^\vee\eta
   +v\otimes L_\epsilon^\vee\eta+1\otimes d_E^\vee\eta
\end{equation}
を左 $D$-線型に延長する。例えば $a_i^2$ に双対な成分では、
二次積が $x_i^2-y_iz_i+z_iy_i=\alpha_1^{(i)}v$ を与え、
$d_E^\vee$ の $-\alpha_1^{(i)}\epsilon^\vee$ がこれを相殺する。
一般にも式\eqref{cmo:eq:curved-ambient-dual}の Leibniz 則により $d_K^2=0$ となる。
PBW filtration に関する associated graded は
$S_0\otimes S_0\otimes\C[v]$ の通常の Koszul 解消である。
各内部次数では filtration が有限なので、この比較から $K\to\C$ は
長さ七の自由解消となる。総階数は $128$ である。

本稿の右加群規約へ戻すときは、
\begin{equation}\label{cmo:eq:curved-right-dual}
 K_i^{\mathrm r}=\underline{\operatorname{Hom}}_D(K_{7-i},D)(-8)
\end{equation}
を用いる。命題\ref{cmo:prop:invariant-weighted-algebra}の
Calabi--Yau 次元七・parameter 八から、これは右剰余体の自由解消である。
以下の行列を転置し、この添字と内部 shift を用いれば、右加群の構成を得る。
非可換係数の積の順序を変えずに通常の転置を同じ左加群上で行う操作とは区別する。

\begin{theorem}[固定値の実行列分解と有限曲率縮約]
\label{cmo:thm:finite-curved-invariant-model}
指定係数 $\alpha^{(1)}=(1,4,9)$、$\alpha^{(2)}=(16,25,49)$、
$w=\Omega_1+100\Omega_2$ に対し、次の構成がある。
\begin{enumerate}
 \item $K_i\to K_{i+2n-1}$、内部次数 $3n$ の有理係数行列
 $\sigma_n$、$0\le n\le4$ を有限個の線型方程式から構成できる。
 $\sigma_0=d_K$ として
 \begin{equation}\label{cmo:eq:curved-shamash-identities}
  \sum_{a+b=n}\sigma_a\sigma_b=
  \begin{cases}w\operatorname{id}_K&n=1,\\0&n\ne1\end{cases}
 \end{equation}
 を満たす。これらを周期化すると、階数 $64+64$ の次数付き行列分解
 $Q_{128}$ を得る。これは剰余体 $\C$ の安定化を表す。
 \item 定数部分 $\partial=d_E^\vee$ の可縮成分を、曲率を保って有限式で除ける。
 得られる行列分解を $Q_{80}$、その微分を $d_{80}$ とすると、階数は $40+40$ であり、
 実際の行列 $I,P,H$ が
 \begin{equation}\label{cmo:eq:curved-contraction-identities}
 \begin{gathered}
 d_{80}^2=w\operatorname{id},\qquad d_{128}I=I\,d_{80},\qquad
 Pd_{128}=d_{80}P,\\
 PI=\operatorname{id},\qquad
 IP=\operatorname{id}-d_{128}H-Hd_{128}
 \end{gathered}
 を満たす。従って $Q_{128}\simeq Q_{80}$ は行列分解の実 homotopy 同値である。
 \item 右加群版の $Q_{80}$ を $\kappa_R$ と書けば、
 \begin{equation}\label{cmo:eq:curved-invariant-generator}
  G_R=\kappa_R\oplus\kappa_R(1)\oplus\kappa_R(2)
 \end{equation}
 は $\mathcal K_R$ を thick 生成する。
 \end{enumerate}
\end{theorem}

\begin{proof}
まず高次 homotopy を構成する。
中心元 $w$ は $H_0(K)=\C$ を零とするので、$w\operatorname{id}_K$ は null-homotopic である。
したがって $d_K\sigma_1+\sigma_1d_K=w\operatorname{id}$ を満たす
$\sigma_1$ がある。一般に、既に低次の $\sigma_a$ を構成したとして
\[
 d_K\sigma_n+\sigma_nd_K=-\sum_{a=1}^{n-1}\sigma_a\sigma_{n-a}
 \qquad(n\ge2)
\]
を順に解く。右辺は既知の恒等式と $w$ の中心性により cycle となり、
$K$ の正の次数での完全性から、基底列ごとの境界持上げで解ける。
この際、既に構成した低い homological 次数の列を固定して進む。
内部次数 $3n$ の部分だけで解けば、全ての行列は内部次数を保つ。
$K_i=0$、$i\notin[0,7]$ なので $n\ge5$ の行列も残る恒等式も自動的に零となる。
これは既知の Shamash 構成
\cite[Definition~3.4.1, Proposition~3.4.2]{CMOEisenbudPeeva}の中心元の場合である。
上の持上げは中心元の乗法以外の可換性を用いず、非可換 $D$ 上で成立する。

各 $K_i$ の基底元の内部重みを $\operatorname{wt}(e)$ とすると、
$\sigma_n$ の $(e,f)$ 成分の係数は
$D_{3n+\operatorname{wt}(e)-\operatorname{wt}(f)}$ に属する。
この有限次元空間は、同じ内部次数の語関係の有理行列から得られる。
従って上の帰納法は、有限個の有理行列の消去として実行可能である。

周期化は
\[
 Q_0=\bigoplus_{p=0}^3K_{2p}(3p),\qquad
 Q_1=\bigoplus_{p=0}^3K_{2p+1}(3p)
\]
とし、$\sum_n\sigma_n$ を $Q_1\to Q_0\to Q_1(3)$ の二つの写像に分ける。
式\eqref{cmo:eq:curved-shamash-identities}が両合成 $w$ を与える。
さらに $R=D/(w)$ 上の Shamash 複体
\[
 T_m=\bigoplus_{j\ge0}(R\otimes_D K_{m-2j})(-3j)
\]
を、$\sigma_a$ が $j$ を $j-a$ へ送る規約で作る。
内部次数の shift により、これらの写像は次数零になる。
$R\otimes_D K$ の homology は $\C$ と $\C(-3)[1]$ であり、
この二つの列を結ぶ $\sigma_1$ の誘導写像は
$d_K\sigma_1+\sigma_1d_K=w$ による connecting isomorphism である。
したがって $j$ の filtration の spectral sequence は正の homological 次数で消え、
$T$ は $\C$ の自由解消となる。
この議論も $w$ の中心正則性だけを用いる。
$T_9=Q_1(-12)$、$T_8=Q_0(-12)$ なので、
$\operatorname{coker}(Q_1\to Q_0)\simeq\Omega_R^8\C(12)$ である。
特異圏の指定周期 $(3)\simeq[2]$ により、これは $\C$ の安定化を表す。

次に定数微分 $\partial$ を縮約する。
一ブロックで $q_i=-\sum_j\alpha_j^{(i)}\xi_{ij}^2$ を掛ける写像は、
次数 $0\to2$ と $1\to3$ のそれぞれで階数一であり、他は零である。
従って、$q_i$ の作用を持つ $E_0$ は、始点の次数が零と一の長さ二の鎖を各一本、
次数一と二の長さ一の鎖を各二本持つ。
$q=q_1\otimes1+1\otimes q_2$ について、二本の長さ二の鎖の tensor 積は
長さ三と一の鎖に分かれる。長さ二と一の tensor 積も合わせると、
\begin{equation}\label{cmo:eq:curved-q-ranks}
 \operatorname{rank}(q:F_r\to F_{r+2})=(1,6,10,6,1)
 \quad (r=0,\ldots,4)
\end{equation}
となる。従って $\partial$ の階数は $24$ で、
$H=H(\C\otimes_D K,\partial)$ の次元は $80$ である。
その homological 次数別次元は $(1,6,14,19,19,14,6,1)$ となり、
偶数部分と奇数部分は各四十次元である。

有限次元の有理線型代数により、両次数を保つ contraction
\[
 (H,0)\xrightarrow{i}(E^\vee,\partial)
 \xrightarrow{p}(H,0),\quad
 pi=1,\quad\partial h_0+h_0\partial=1-ip,
 \quad h_0^2=ph_0=h_0i=0
\]
を取る。$\partial$ は $F^\vee$ から $(F\epsilon)^\vee$ へ移すので、
$h_0$ は $\epsilon$ 成分を一から零へ減らし、元の homological 次数を一上げるように選べる。
さらに $h_0$ の非零な始点の homological 次数は一から五までである。

$\Delta=d_{128}-\partial$ と置く。
$d_K-\partial$ の項を $h_0$ の後で使うと、$\epsilon$ 成分は零のままとなり、
それより左に現れる次の $h_0$ は零になる。
一方、$\sigma_nh_0$、$n\ge1$ は元の homological 次数を $2n$ 上げる。
従って $(\Delta h_0)^4$ が非零なら、右から三つの段階でこの次数が少なくとも六増え、
第四の $h_0$ の始点が七以上になる。しかしその始点は高々五なので矛盾する。
こうして
\begin{equation}\label{cmo:eq:curved-finite-inverse}
 (\Delta h_0)^4=0,\qquad
 B=(1+\Delta h_0)^{-1}
   =1-\Delta h_0+(\Delta h_0)^2-(\Delta h_0)^3
\end{equation}
を得る。この逆行列は $D$ 上の有限式であり、形式級数の収束を仮定しない。

$A=B\Delta$ と置き、
\begin{equation}\label{cmo:eq:curved-transfer-formulas}
 d_{80}=pAi,\quad I=i-h_0Ai,\quad P=p-pAh_0,\quad
 H=h_0-h_0Ah_0
\end{equation}
と定める。これらは全て実際の有理係数行列である。
恒等式 $\partial h_0+h_0\partial=1-ip$ を使うと
\[
 (1+\Delta h_0)
 (\partial A+A\partial+AipA)
 (1+h_0\Delta)=w\operatorname{id}
\]
となる。$w$ が中心で $h_0^2=ph_0=h_0i=0$ であることから、
この式と $A+\Delta h_0A=\Delta=A+Ah_0\Delta$ を代入して
式\eqref{cmo:eq:curved-contraction-identities}を得る。
特に転送後も曲率は厳密に $w$ であり、単に曲率を忘れた複体の縮約ではない。
式\eqref{cmo:eq:curved-right-dual}を取ると、転置した $P,I,H$ が右加群版の
包含・射影・homotopy を与える。

最後に生成性を示す。$C=A_0\rtimes\mu_2$ とし、
次数付き特異圏上の関手
\[
 j_{\mathrm{sg}}=(-)\otimes_R e_+C,\qquad
 j_{\mathrm{sg}}^*=(-)e_+
\]
を取る。$e_+C$ は左 $R$-自由・右 $C$-射影、$Ce_+$ は右 $R$-自由なので、
両関手は perfect 複体を保ち、特異圏へ降りる。
また $j_{\mathrm{sg}}^*j_{\mathrm{sg}}=\operatorname{id}$ である。
\S\ref{mix:sec:equivariant-window}の三対象 window と有限群同変化により、
$D_{\mathrm{sg}}^{\mathbb Z}(C)$ は
$\C_C^\pm(r)$、$r=0,1,2$ により thick 生成される。
ここで $j_{\mathrm{sg}}^*\C_C^+(r)=\C_R(r)$、
$j_{\mathrm{sg}}^*\C_C^-(r)=0$ である。
従って $D_{\mathrm{sg}}^{\mathbb Z}(R)$ は三つの $\C_R(r)$ で生成され、
命題\ref{cmo:prop:invariant-weighted-algebra}の MF 比較が
式\eqref{cmo:eq:curved-invariant-generator}を与える。
この議論は、異なる qgr window での $j$ の同定を追加の仮定として用いない。
\end{proof}

\paragraph{行列の再現と接着への入力.}
\texttt{computations/verify\_mixed\_curved\_dual.py} は、二次関係の直交空間から
$E_0$ を作り、各内部次数の語関係の有理消去で係数環を計算する。
境界持上げの行列は、基底と pivot の順序を固定して生成する。
指定入力で得た $\sigma_1,\sigma_2$ の係数単項式数はそれぞれ $1312,624$ で、
この選択では $\sigma_3=\sigma_4=0$ となる。
また $(\Delta h_0)^3=0$ も検算されるが、定理の有限性には四乗零の一般評価を用いた。
転送微分の係数単項式数は $8228$ である。
検算器は全ての Shamash 恒等式、内部次数、転送後の曲率、
$I,P,H$ の五つの比較恒等式を有理数上で直接確認する。
\texttt{--export} を指定すれば、後続の比較に用いる全行列を出力できる。
これらの有限恒等式の実行確認と、解消の完全性・生成性の上の証明を区別する。

\begin{corollary}[実 dg 模型と制限核]
\label{cmo:cor:curved-invariant-end-model}
$\Gamma_R=\operatorname{End}_{\mathrm{MF}^{\mathbb Z}(D,w)}(G_R)$ と置くと、
\[
 \mathcal K_R\xrightarrow{\ \operatorname{RHom}(G_R,-)\ }
 \operatorname{Perf}(\Gamma_R)
\]
は実 dg Morita 同値であり、逆関手は $-\otimes_{\Gamma_R}^{\mathbb L}G_R$ である。
$\Gamma_R$ の微分と積は、それぞれ行列分解の supercommutator と行列合成で与えられる。
その各固定 cohomological 次数・内部次数の係数計算は有限である。
また球面制限 $\mathcal L:\mathcal K_R\to\mathcal K_S$ に対する
$\mathcal L(G_R)$ は、上の行列の係数を $v=0$ とすることで表される。
\end{corollary}

\begin{proof}
前定理の thick 生成性から dg Yoneda--Morita の定理を適用する。
微分の二乗は $wf-fw=0$ であり、全行列に内部 shift を付けて Hom 複体を作れる。
$D/(v)$ 上への制限は有限自由行列分解の係数の底変更で計算でき、
$I,P,H$ の恒等式も底変更後に成立する。
従ってこの表示は、既存の球面制限の実核に対する生成対象上の具体的な表示を与える。
\end{proof}

\paragraph{可換三様体との比較に残る積の問題.}
この有限縮約は、$(v,L)$ による Kn\"orrer 同値ではない。
特に $F$ の積を変えずに $\epsilon$ と各次数一生成元の反交換子だけを
可逆な中心周期変数 $t$ の定数倍に置き換える方法は使えない。
実際、結合的なそのような拡大で
$\epsilon a+a\epsilon=t\ell_a$、
$\epsilon b+b\epsilon=t\ell_b$、
$\epsilon c+c\epsilon=t\ell_c$ とすれば、
$t^{-1}[\epsilon,-]_{\mathrm{super}}$ は $E_0$ 上の奇導分になる。
関係 $bc+a^2=0$ の像は $\ell_b c-\ell_c b$ だから
$\ell_b=\ell_c=0$ であり、巡回的に $\ell_a=0$ も従う。
指定 $L$ の係数は非零なので、この厳密結合的な ansatz は成立しない。
実際の $A_\infty$ 模型では、積の変更または高次積による補正が必要である。

可換三様体 $Y$ との実比較と制限の intertwiner は、
定理\ref{cmo:thm:fixed-zero-commutative-kernel}で球面境界の方法により構成する。
その核をこの $\Gamma_R$ と可換側の実生成代数 $\Gamma_Y$ の間の
行列双加群としてさらに具体化することは、別の計算課題である。
系\ref{cmo:cor:curved-invariant-end-model}はその source と制限を
固定値で計算できる形にする。曲率付き高次積の全式や直接の Morita 消去を、
総階数八十という数値だけから結論しない。

% !TEX root = ../../mixed-sklyanin-cubics.tex
% Current result: 2026-11-22 / v53.
\subsection{球面制限の twist と有限生成}
\label{cmo:subsec:spherical-comparison}

命題\ref{cmo:prop:root-order-comparison-criterion}では、
不変環の残余圏の比較を、最初から同値として構成することを仮定した。
ここでは制限関手の周期を利用して、この要求を弱める。
実際の候補関手と制限の intertwiner が与えられたとき、
その完全忠実性と本質的全射性を球面随伴の三角から検査できる。
以下の cone と自然変換は全て perfect 双加群核の導来圏で取る。
単に各対象に同型が存在するという条件ではなく、
合成と両立する dg 自然変換を用いる。

\begin{proposition}[制限の twist と二つの周期]
\label{cmo:prop:spherical-restriction-twist}
命題\ref{cmo:prop:invariant-restriction-spherical}の
$F=\mathcal L|_{\mathcal K_R}:\mathcal K_R\to\mathcal K_S$ に対して
\[
 \mathsf T_F=\operatorname{Cone}
   (FF^R\xrightarrow{\epsilon}\operatorname{id}_{\mathcal K_S}),
 \qquad
 \mathsf C_F=\operatorname{Cone}
   (\operatorname{id}_{\mathcal K_R}\xrightarrow{\eta}F^RF)
\]
と置く。このとき、内部 shift に対応する境界の rotation $\varrho_S$ を用いて
\begin{equation}
\label{cmo:eq:spherical-restriction-periods}
 \mathsf T_F\simeq\varrho_S,\qquad
 \mathsf T_F^3\simeq[2],\qquad
 \mathsf C_F^3\simeq[-1]
\end{equation}
である。同じ結論は可換不変環 $R^{\mathrm{cl}}$ とその境界にも成立する。
\end{proposition}

\begin{proof}
前と同じく $a:\mathcal K_R\hookrightarrow\mathcal T_R$、
$a\dashv p_R$、$\rho:\mathcal T_S\to\mathcal T_R$ と書く。
$\mathcal T_R$ の Serre 関手が $(-5)[5]$ なので、
五本の例外線束を巡回移動すると
\[
 \mathcal T_R=
 \langle\OO_R(-5),\OO_R(-4),\OO_R(-3),
                  \OO_R(-2),\OO_R(-1),\mathcal K_R\rangle
\]
となる。従って $a p_R$ は、この負の五本の線束に関する right mutation である。
$N\in\mathcal K_S$ に対して、中心元 $v$ の長さ一の解消とその双対から
\begin{equation}
\label{cmo:eq:spherical-twist-line-vanishing}
 R\operatorname{Hom}_{\mathcal T_R}(\rho N,\OO_R(r))
 \simeq
 R\operatorname{Hom}_{\mathcal T_S}(N,\OO_S(r+2))[-1]
\end{equation}
を得る。$r=-5,-4,-3$ では右辺は零である。
実際、$\mathcal T_S$ の Serre 関手は $(-3)[4]$ なので、
この消滅は $\OO_S(0),\OO_S(1),\OO_S(2)$ から $N$ への消滅と同値である。
right mutation を表示した順に行うと、残る correction は最後の二本だけである。
従って自然な三角
\[
 a p_R\rho N\longrightarrow\rho N\longrightarrow E_N,
 \qquad E_N\in\langle\OO_R(-2),\OO_R(-1)\rangle
\]
が得られる。

全圏の制限の余単位 $\mathcal L\rho N\to N$ は、
$v$ の二項解消によって
\[
 \operatorname{Cone}(\mathcal L\rho N\to N)=N(-2)[2]
\]
を満たす。上の三角へ $\mathcal L$ を適用し、この余単位と合成する。
得られる $F F^R N\to N$ は残余関手の余単位そのものであり、
octahedron により
\begin{equation}
\label{cmo:eq:spherical-twist-mutation-triangle}
 \mathsf T_F N\longrightarrow N(-2)[2]
                  \longrightarrow\mathcal L E_N[1]
\end{equation}
を得る。
第一項は $\mathcal K_S$ に属し、第三項は
$\langle\OO_S(-2),\OO_S(-1)\rangle$ に属する。
Serre 双対性により $\mathcal K_S$ はこの二本の線束の左直交に含まれる。
従って\eqref{cmo:eq:spherical-twist-mutation-triangle}は
right mutation を特徴づける三角であり、
\[
 \mathsf T_F\simeq
 R_{\langle\OO_S(-2),\OO_S(-1)\rangle}\circ(-2)[2]
 \simeq\varrho_S^{-2}[2]\simeq\varrho_S.
\]
最後の同型は三次行列分解の周期 $\varrho_S^3\simeq[2]$ による。
同じ周期と\eqref{cmo:eq:invariant-restriction-cotwist}から
\[
 \mathsf C_F^3\simeq\varrho_R^{-6}[3]\simeq[-1]
\]
を得る。
この証明に用いたのは parameter 五と三、中心正則元の次数二、三次周期、
および表示した解消・随伴だけなので、可換側にも同じ証明が適用できる。
\end{proof}

\begin{lemma}[周期による conservativity と split-generation]
\label{cmo:lem:periodic-spherical-generation}
$\mathcal A,\mathcal B$ を体 $k$ 上の proper な冪等完備 dg 圏とし、
$F:\mathcal A\to\mathcal B$ を右随伴 $R_F$ を持つ dg 関手とする。
単位と余単位の cone が自己同値 $\mathsf C,\mathsf T$ であり、
\[
 \mathsf C^3\simeq[-1],\qquad\mathsf T^3\simeq[2]
\]
を満たすと仮定する。このとき $F$ は conservative であり、
\begin{equation}
\label{cmo:eq:spherical-two-generations}
 \mathcal B=\operatorname{thick}(F\mathcal A),
 \qquad
 \mathcal A=\operatorname{thick}(R_F\mathcal B)
\end{equation}
である。特に $G$ が $\mathcal B$ の split-generator なら、
$R_FG$ は $\mathcal A$ の split-generator である。
\end{lemma}

\begin{proof}
$FM=0$ なら単位三角から $\mathsf C M\simeq M[1]$ となる。
三回適用すると $M[-1]\simeq M[3]$、従って $M\simeq M[4]$ である。
proper 性により $R\operatorname{Hom}(M,M)$ は有界な有限次元複体である。
非零の恒等写像が任意に大きい次数へ移ることはできないため $M=0$ となる。

次に、余単位三角を $N,\mathsf T N,\ldots,\mathsf T^{n-1}N$ に適用する。
それぞれの射 $\mathsf T^jN\to\mathsf T^{j+1}N$ の cone は
$F R_F(\mathsf T^jN)[1]$ である。
従って有限回の octahedron により、合成
\[
 N\longrightarrow\mathsf T^nN
\]
の cone は $\operatorname{thick}(F\mathcal A)$ に入る。
$n=3m$ とし $m$ を十分大きく取れば、proper 性から
$\operatorname{Hom}(N,N[2m])=0$ であり、この合成は零になる。
その cone は $N[2m]\oplus N[1]$ なので、$N$ は同じ thick 部分圏に属する。

最後に $\mathsf H=\mathsf C[-1]$ と置く。
単位の fiber 三角は
\[
 \mathsf H M\longrightarrow M\longrightarrow R_F FM
\]
であり、$\mathsf H^3\simeq[-4]$ である。
この三角を $\mathsf H^jM$ に反復すると、
合成 $\mathsf H^nM\to M$ の cone は
$\operatorname{thick}(R_F\mathcal B)$ に入る。
$n=3m$、$m$ 十分大では
$\operatorname{Hom}(M[-4m],M)=0$ なので合成は零であり、
その cone の直和因子 $M$ も同じ thick 部分圏に属する。
$R_F$ は exact で直和因子を保つため、最後の生成対象についての主張も従う。
\end{proof}

この補題は有限個の対象の Hom の有界性を用いる。
Verdier 商の proper 性や $\operatorname{Ind}(F)$ の conservativity を仮定せず、
それらを結論してもいない。

\subsection{候補核の完全忠実性を強制する周期}

比較射を正確に固定するため、次の一般の記号を用いる。
$F:\mathcal A\to\mathcal B$、$F':\mathcal A'\to\mathcal B'$ の右随伴を
$R_F,R_{F'}$ とし、単位・余単位を $\eta,\epsilon$ と $\eta',\epsilon'$、
unshifted cotwist と twist を $\mathsf C,\mathsf T$ と
$\mathsf C',\mathsf T'$ と書く。
perfect 核を持つ関手
\[
 P:\mathcal A\longrightarrow\mathcal A',\qquad
 E:\mathcal B\xrightarrow{\sim}\mathcal B',\qquad
 \psi:EF\xrightarrow{\sim}F'P
\]
が与えられたとする。ここで $P$ はまだ同値と仮定しない。
$\psi$ の随伴 mate を
\begin{equation}
\label{cmo:eq:spherical-canonical-mate}
 \alpha:P R_F\longrightarrow R_{F'}E
\end{equation}
とする。具体的にはこれは
\[
 F'P R_F\xrightarrow{\psi_{R_F}^{-1}}EF R_F
                   \xrightarrow{E\epsilon}E
\]
を $F'\dashv R_{F'}$ により transpose した射である。
随伴の単位の自然性から、
\[
 P\eta:P\longrightarrow P R_F F,\qquad
 \eta'P:P\longrightarrow R_{F'}F'P
\]
の間には、第一項が恒等、第二項が
$R_{F'}\psi\circ\alpha_F$ である可換な射がある。
この射の cone を
\begin{equation}
\label{cmo:eq:spherical-canonical-cotwist-map}
 c_P:P\mathsf C\longrightarrow\mathsf C'P
\end{equation}
と定める。
以下で必要なのはこの標準射であり、二つの関手の抽象的な同型だけではない。

\begin{theorem}[cotwist と Hom の欠陥による比較基準]
\label{cmo:thm:spherical-cotwist-comparison}
$\mathcal A,\mathcal A',\mathcal B,\mathcal B'$ は smooth proper dg 圏で、
$F,F'$ は perfect 核を持つ spherical 関手とする。
さらに
\begin{equation}
\label{cmo:eq:spherical-comparison-period-input}
 \mathsf C^3\simeq[-1],\qquad
 (\mathsf C')^3\simeq[-1],\qquad
 \mathsf T^3\simeq[2]
\end{equation}
と仮定する。上記の $P,E,\psi$ に対して
標準射 $c_P$ が同型なら、$P$ は dg Morita 同値である。
\end{theorem}

\begin{proof}
$M,N\in\mathcal A$ に対し、完全忠実性の欠陥を
\[
 D(M,N)=\operatorname{Cone}\left(
 R\operatorname{Hom}_{\mathcal A}(M,N)
   \longrightarrow R\operatorname{Hom}_{\mathcal A'}(PM,PN)\right)
\]
と置く。両圏が proper なので、これは $k$ 上の perfect 複体である。
単位三角へ $R\operatorname{Hom}(M,-)$ を適用し、随伴を用いると
\[
 R\operatorname{Hom}_{\mathcal A}(M,N)
 \longrightarrow R\operatorname{Hom}_{\mathcal B}(FM,FN)
 \longrightarrow R\operatorname{Hom}_{\mathcal A}(M,\mathsf C N)
\]
という三角を得る。$F'$ の三角にも同じ操作を行う。
これらの三角の間の第一の射は $P$ による射、
第二の射は $E$ と $\psi$ による同型、
第三の射は $P$ と標準射 $c_P$ による射である。
cone の三角と $c_P$ の可逆性から
\begin{equation}
\label{cmo:eq:spherical-hom-defect-period}
 D(M,\mathsf C N)\simeq D(M,N)[1]
\end{equation}
となる。これを三回適用し、$\mathsf C^3\simeq[-1]$ を使うと
\[
 D(M,N)[-1]\simeq D(M,N)[3]
\]
を得る。有界な有限次元複体は、非零ならこの四周期性を持てない。
従って $D(M,N)=0$ であり、$P$ は完全忠実である。

perfect 核の関手は smooth proper 圏の間で左右随伴を持つので、
像 $\mathcal I=P\mathcal A$ は admissible である。
$c_P$ と $\mathsf C$ の可逆性から
$\mathsf C'\mathcal I=\mathcal I$ である。
$N'\in\mathcal I^\perp$ とすると、任意の $M\in\mathcal A$ について
\[
 R\operatorname{Hom}_{\mathcal A'}(PM,N')=0,
 \qquad
 R\operatorname{Hom}_{\mathcal A'}(PM,\mathsf C'N')=0.
\]
第二の消滅では $(\mathsf C')^{-1}\mathcal I=\mathcal I$ を用いた。
$F'$ の単位三角により
\[
 R\operatorname{Hom}_{\mathcal B'}(F'PM,F'N')=0
\]
となる。補題\ref{cmo:lem:periodic-spherical-generation}の生成性と
$F'P\simeq EF$ により、$F'P\mathcal A$ は $\mathcal B'$ を split-generate する。
従って $F'N'=0$ である。
同補題の conservativity の証明を $F'$ と $\mathsf C'$ に適用すると $N'=0$ を得る。
従って $\mathcal I^\perp=0$、つまり $P$ は本質的全射でもある。
最初から perfect dg 核を用いているので、得られた同値は dg Morita 同値である。
\end{proof}

この証明は Hochschild homology の一致や数値格子の相殺を用いない。
完全忠実性に失敗する Hom 複体自体に四周期性を課し、proper 性でそれを排除している。

\subsection{twist の標準射による比較と可換境界}

式\eqref{cmo:eq:spherical-canonical-mate}を余単位と組み合わせると、
\[
 \epsilon' E\circ F'\alpha\circ\psi_{R_F}=E\epsilon
\]
である。従って余単位の cone の間に標準射
\begin{equation}
\label{cmo:eq:spherical-canonical-twist-map}
 t_P:E\mathsf T\longrightarrow\mathsf T'E
\end{equation}
が定まり、twist 三角の射 $\nu:\operatorname{id}\to\mathsf T$、
$\nu':\operatorname{id}\to\mathsf T'$ について
\begin{equation}
\label{cmo:eq:spherical-twist-unit-compatibility}
 t_P\circ E\nu=\nu'E
\end{equation}
が成立する。

\begin{corollary}[twist mate による比較基準]
\label{cmo:cor:spherical-twist-comparison}
定理\ref{cmo:thm:spherical-cotwist-comparison}の周期の仮定の下で、
標準射 $t_P$ が同型なら $P$ は dg Morita 同値である。
さらに $k=\C$、$HH^0(\mathcal B)=\C$ とし、
何らかの dg 自然同型
\[
 \beta:E\mathsf T\xrightarrow{\sim}\mathsf T'E
\]
が存在すると仮定する。$\nu'\ne0$ なら $t_P$ は自動的に同型である。
この場合、候補 $P$ に対して要求する比較データは
実際の制限 intertwiner $\psi$ だけでよい。
\end{corollary}

\begin{proof}
$t_P$ が同型なら、余単位三角の間の射において
第二項 $E\to E$ と第三項 $t_P$ が同型である。
従って第一項の射 $F'\alpha\circ\psi_{R_F}$、従って $F'\alpha$ も同型である。
$F'$ は補題\ref{cmo:lem:periodic-spherical-generation}の証明により conservative なので、
$\alpha$ は同型である。
単位三角の間の射に戻ると、第一項と第二項が同型であるため
$c_P$ も同型となり、定理を適用できる。

追加の仮定の下では、$E\mathsf T$ は同値なので
\[
 H^0\operatorname{RHom}(E\mathsf T,\mathsf T'E)
 \xrightarrow[\ \beta^{-1}\circ(-)\ ]{\ \sim\ }
 H^0\operatorname{REnd}(E\mathsf T)
 \simeq HH^0(\mathcal B)=\C.
\]
従って $t_P$ は $\beta$ の定数倍である。
式\eqref{cmo:eq:spherical-twist-unit-compatibility}と $\nu'\ne0$ により
この定数は零でないので、$t_P$ は同型である。
\end{proof}

本例に適用すると
\[
 (\mathcal A,\mathcal B,F)=(\mathcal K_R,\mathcal K_S,\mathcal L),
 \qquad
 (\mathcal A',\mathcal B',F')=
 (\mathcal K_{R^{\mathrm{cl}}},\mathcal K_{S^{\mathrm{cl}}},\mathcal L^{\mathrm{cl}})
\]
であり、周期の仮定は命題\ref{cmo:prop:spherical-restriction-twist}により満たされる。
可換境界 $\mathcal B'=\Ku(Z)$ は滑らかな三次四様体の残余圏で、
$HH^0(\mathcal B')=\C$ である
\paperref{main}{prop:nf-hh}。
従って境界比較 $E$ が rotation を保ち、可換側の twist 射 $\nu'$ が非零なら、
任意の perfect 候補 $P$ と実際の
\[
 E\mathcal L\xrightarrow{\sim}\mathcal L^{\mathrm{cl}}P
\]
から $P$ の同値性が従う。
これを命題\ref{cmo:prop:root-order-comparison-criterion}へ代入すれば、
同変五様体の貼合せ核と通常の五様体の比較核を構成できる。
この最後の簡略版では、境界比較の rotation 両立と $\nu'\ne0$ を明示的な入力とする。
これらは補題\ref{cmo:lem:heisenberg-boundary-comparison}と
補題\ref{cmo:lem:spherical-bockstein}で確認する。
この判定は候補 $P$ を入力に取る。
次に、twist 射の比較から $P$ と制限 intertwiner 自体を構成する。

\subsection{実 comonad からの有限 Morita 模型}

次の記述は、候補関手を構成する別の具体的な目標を与える。
twist の同型類だけから源の圏が決まるという主張ではない。
twist とその section から Kleisli 圏を用いて球面関手を復元する考え方は
Segal \cite[Section~3.2]{CMOSegalSphericalTwists}で論じられている。
同節では dg coherence の問題も明示されている。
さらに Christ \cite[Section~4.1, Example~4.10]{CMOChristSphericalMonadic}は、
underlying functor と unit が同じでも乗法の異なる monad を用いて、
一般には twist と section だけから復元できないことを示した。
以下では、右随伴像の有限生成を周期から確認した上で、
三次四様体に固有の全 arity の消滅を別に証明して比較する。

\begin{proposition}[右随伴の生成対象と coKleisli 代数]
\label{cmo:prop:spherical-comonad-model}
補題\ref{cmo:lem:periodic-spherical-generation}の仮定の下で、
$\mathcal A,\mathcal B$ は smooth proper であるとする。
$G$ を $\mathcal B$ の split-generator とし、
随伴から来る実 comonad を
\[
 K=F R_F,\qquad
 \epsilon:K\longrightarrow\operatorname{id}_{\mathcal B},\qquad
 \Delta=F\eta R_F:K\longrightarrow K^2
\]
と書く。このとき
\begin{equation}
\label{cmo:eq:spherical-comonad-morita-model}
 \mathcal A\simeq\operatorname{Perf}(\Gamma),\qquad
 \Gamma=R\operatorname{End}_{\mathcal A}(R_FG)
       \simeq R\operatorname{Hom}_{\mathcal B}(KG,G).
\end{equation}
右辺の積は、coKleisli 表示では
\begin{equation}
\label{cmo:eq:spherical-comonad-product}
 g\star f=g\circ K(f)\circ\Delta_G,\qquad 1=\epsilon_G
\end{equation}
である。dg 模型に高次の構造写像が必要な場合には、
実 adjunction の homotopy-coherent comonad 構造を含めて読む。

特に $E:\mathcal B\simeq\mathcal B'$ が、余単位・余乗法と
その全 coherence を含む comonad 同型 $EK\simeq K'E$ を与えるなら、
式\eqref{cmo:eq:spherical-comonad-morita-model}の代数の比較から
実際の perfect 核 $P:\mathcal A\simeq\mathcal A'$ と
$EF\simeq F'P$ を構成できる。
\end{proposition}

\begin{proof}
補題\ref{cmo:lem:periodic-spherical-generation}により $R_FG$ は生成対象である。
従って $R\operatorname{Hom}_{\mathcal A}(R_FG,-)$ と
生成対象との導来 tensor は dg Morita 同値を与える。
これは stable Kleisli 構成
\cite[Proposition~1.11]{CMOChristSphericalMonadic}の双対と冪等完備化を、
一つの生成対象の自己準同型代数で書いたものである。
随伴の同型
\[
 R\operatorname{Hom}_{\mathcal A}(R_FN,R_FN')
 \simeq R\operatorname{Hom}_{\mathcal B}(F R_FN,N')
\]
へ単位・余単位の恒等式を適用すると、合成は
\eqref{cmo:eq:spherical-comonad-product}となる。
これは coKleisli 圏の合成そのものであり、
実際の随伴の構造を使うので合成の結合則も保持する。
具体的には $f:KG\to G$ に対応する自己準同型は
$R_F(f)\circ\eta_{R_FG}$ であり、その $F$ による像は
\[
 F\bigl(R_F(f)\circ\eta_{R_FG}\bigr)
   =K(f)\circ\Delta_G
\]
である。この式は生成対象上の $F$ の作用も comonad の構造から決まることを示す。

comonad 同型が与えられれば、$G'=EG$ として
\[
 R\operatorname{End}_{\mathcal A}(R_FG)
 \simeq R\operatorname{End}_{\mathcal A'}(R_{F'}G')
\]
を積と高次構造を含めて得る。
この代数の比較と二つの生成対象を tensor すれば $P$ の perfect 核が得られる。
より具体的には、$\Gamma'$ を右辺の代数とし、その $A_\infty$ 比較の graph 双加群を
bar 解消により dg 双加群にすると、$\Gamma'$ を underlying right module とする
Morita 核を得る。二つの生成対象の Morita 核と合成して $P$ とする。
両自己準同型代数は smooth proper であり、graph 双加群は underlying
$k$-複体として perfect なので、$\Gamma^{\mathrm{op}}\otimes\Gamma'$ 上も perfect である。
実際、smooth な対角の有限解消を tensor すれば、任意の $k$-perfect 加群は perfect になる。
従って bar 模型が無限であることは、得られた導来核の perfect 性を妨げない。
生成対象上の同型
\[
 EF(R_FG)=EKG\simeq K'EG=F'(R_{F'}G')
\]
は、comonad の比較により生成対象の自己準同型代数の作用と両立する。
従って twisted complex と直和因子へ延長して $EF\simeq F'P$ を得る。
\end{proof}

ここで $K$ の underlying functor は
\[
 K\simeq\operatorname{Cone}
     (\operatorname{id}_{\mathcal K_S}\xrightarrow{\nu}\varrho_S)[-1]
\]
と書ける。一般には、この三角だけから余乗法と全 coherence の一致は従わない。
三次四様体の境界では、次の消滅がその追加構造を制御する。

\begin{lemma}[正規化した余乗法の一意性]
\label{cmo:lem:normalized-comonad-uniqueness}
合成 tensor が各変数について exact な dg 核の monoidal 圏で、
余単位の射 $\epsilon:K\to\mathbf1$ を固定し、
$J=\operatorname{fib}(\epsilon)$ と置く。
各 $n\ge2$ と $q\le0$ について
\begin{equation}
\label{cmo:eq:normalized-comonad-vanishing}
 H^qR\operatorname{Hom}(K,J^{\otimes n})=0
\end{equation}
と仮定する。
この $(K,\epsilon)$ 上に二つの実 counital homotopy-coherent comonad 構造があるなら、
underlying kernel 上で恒等となる同値により両者を同定できる。
より正確には、underlying counit arrow を固定したこの構造の空間は、
非空なら可縮である。
\end{lemma}

\begin{proof}
核の monoidal stable $\infty$-category を $\mathscr M$ とする。
必要なら対象を含む小さい stable tensor 閉部分圏に制限する。
opposite category $\mathscr M^{\mathrm{op}}$ では、
$\epsilon$ は unit $\mathbf1\to K$、$J$ はその cofiber となる。
この小さい stable monoidal 圏を Ind 完備化すれば、
tensor を各変数の colimit と可換に延長できる。
得られる presentably monoidal stable 圏への包含は完全忠実かつ monoidal であり、
元の対象の有限 tensor と mapping spectrum は変わらない。
従って、固定 unit を持つ unital $A_n$-構造の空間を扱う
\cite[Theorem~3.3 and Lemma~3.11]{CMOWillumsgaardAssociativity}を適用できる。

同文献の記述で正規化の意味を確認する。
$n$ 個の出力のそれぞれへ $\epsilon$ を施す全ての射を集めた $n$-cube の
初項は $K^{\otimes n}$ であり、非空部分集合の出力へ $\epsilon$ を施した項は
残る $K$ の tensor である。
exact 性により、この cube の total fiber は $J^{\otimes n}$ である。
従って全ての counit reduction とその homotopy を固定した後の複体は
\begin{equation}
\label{cmo:eq:normalized-comonad-complex}
 N_n=R\operatorname{Hom}(K,J^{\otimes n})
\end{equation}
である。これは unit や counit の chain-level splitting を必要としない。

$(K,\epsilon)$ 上の counital $A_n$-構造の空間を $\mathcal S_n$ と書く。
opposite category での unit を固定した空間を取るため、$\mathcal S_1$ は一点である。
$n$ 番目の associahedron の次元は $n-2$ であり、
上記の定理により $\mathcal S_n\to\mathcal S_{n-1}$ の各非空 fiber は、
既存の lift を基点に選ぶと
\begin{equation}
\label{cmo:eq:normalized-comonad-structure-fiber}
 \Omega^\infty N_n[2-n]
\end{equation}
と同値である。ここで $\mathscr M^{\mathrm{op}}$ の suspension は
$\mathscr M$ の $[-1]$ なので、shift は確かに $[2-n]$ となる。
この空間の $i$ 番目の homotopy group は、$i\ge0$ について
\[
 H^{2-n-i}N_n=0
\]
である。従って全ての非空 fiber は可縮である。
仮定により少なくとも一つの全 comonad 構造が存在するので、
各 $\mathcal S_n$ は非空である。
$\mathcal S_1$ から帰納すると、基底が可縮で一つの非空 fiber が可縮であるため、
各 $\mathcal S_n$ も可縮になる。
全 coherent associative 構造の空間は $\varprojlim_n\mathcal S_n$ であるから、
この極限も可縮である。
その二点を結ぶ path が、underlying counit arrow を固定した comonad 同値を与える。

cofibrant dg 模型でこの path を表示すると、比較の第 $n$ 項は次数 $1-n$、
その未消去部分は $N_n$ の次数 $2-n$ の cocycle となる。
仮定は各段でその cocycle の primitive を選べることを意味する。
この表示は正規化した $A_\infty$ 操作と一致する
\cite[Definition~4.1 and Lemma~4.4]{CMOBurkeStrictUnit}。
双加群の完成 tensor による表示
\cite[Sections~2.1 and~2.3]{CMONegronVolkovWitherspoon}では各 arity は有限和であり、
解析的な収束は要求しない。
ただし全 coherence が得られる根拠は、その形式表示だけでなく、
上で検証した unital structure space の tower である。
\end{proof}

\begin{proposition}[三次四様体境界における全 arity 消滅]
\label{cmo:prop:cubic-boundary-comonad-vanishing}
$Z\subset\mathbb P^5$ を滑らかな三次四様体、
$\mathcal B=\Ku(Z)$、$\varrho$ をその rotation とする。
任意の dg 自然変換 $\nu:\operatorname{id}\to\varrho$ に対して
\[
 K=\operatorname{Cone}(\nu)[-1],\qquad
 \epsilon:K\to\operatorname{id},\qquad
 J=\operatorname{fib}(\epsilon)=\varrho[-1]
\]
と置く。この $(K,\epsilon)$ は
\eqref{cmo:eq:normalized-comonad-vanishing}を満たす。
従って、これに実 counital comonad 構造が存在すれば、
その全高次構造は補題\ref{cmo:lem:normalized-comonad-uniqueness}の意味で一意である。
\end{proposition}

\begin{proof}
方程式を $G$ とする。次数付き行列分解の diagonal sector 公式
\cite[Theorem~5.39]{NFBFKKernels}、
\cite[Theorem~3.5 and Remark~3.6]{NFLZSerre}を用いる。
$H^aR\operatorname{Hom}(\operatorname{id},\varrho^b)$ の恒等 sector は、
$a$ が偶数なら $\operatorname{Jac}(G)_{b+3a/2}$ であり、
$a$ が奇数なら零である。
非恒等な二つの $\mu_3$ sector は六変数全てを動かす。
従ってこの部分が現れるためには
\[
 a\equiv6\pmod2,\qquad 2b+3a=6
\]
が必要である。
この条件は、主稿の\paperref{main}{prop:nf-hh}の計算に内部 shift $b$ を加えたものである。

$J^n=\varrho^n[-n]$ であるから、三角 $K\to\operatorname{id}\to\varrho$ により、
$H^qR\operatorname{Hom}(K,J^n)$ の消滅には
\[
 H^qR\operatorname{Hom}(\operatorname{id},J^n)=0,\qquad
 H^qR\operatorname{Hom}(\varrho[-1],J^n)=0
\]
を示せばよい。
前者の恒等 sector の Jacobian 次数は $(3q-n)/2$、
後者では $(3q+1-n)/2$ である。
$n\ge2$、$q\le0$ のときはいずれも負なので零になる。
非恒等 sector の必要条件も、それぞれ
\[
 3q-n=6,\qquad 3q+1-n=6
\]
となり、この範囲では満たされない。従って必要な全ての次数が消える。
\end{proof}

\begin{corollary}[twist 射を同定した後の実比較核]
\label{cmo:cor:spherical-unit-reconstruction}
本例の二つの球面制限 $F:\mathcal K_R\to\mathcal K_S$、
$F':\mathcal K_{R^{\mathrm{cl}}}\to\mathcal K_{S^{\mathrm{cl}}}$ に対して、
perfect 核を持つ境界同値と rotation の同型
\[
 E:\mathcal K_S\xrightarrow{\sim}\mathcal K_{S^{\mathrm{cl}}},
 \qquad \beta:E\varrho_S\xrightarrow{\sim}\varrho_{S^{\mathrm{cl}}}E
\]
を与える。この同型の下で
\begin{equation}
\label{cmo:eq:spherical-unit-matching}
 \beta\circ E\nu=\nu' E
\end{equation}
が核の導来圏の射として成立すれば、
不変環の比較核 $P:\mathcal K_R\simeq\mathcal K_{R^{\mathrm{cl}}}$ と
$EF\simeq F'P$ を実際に構成できる。
従って、命題\ref{cmo:prop:root-order-comparison-criterion}により
通常の五様体残余圏の dg Morita 同値も構成できる。
\end{corollary}

\begin{proof}
式\eqref{cmo:eq:spherical-unit-matching}の homotopy を選んで fiber を取ると、
underlying kernel と counit の同型 $E K\simeq K'E$ を得る。
両辺には元の随伴から実際の comonad 構造がある。
命題\ref{cmo:prop:cubic-boundary-comonad-vanishing}と
補題\ref{cmo:lem:normalized-comonad-uniqueness}により、
この同型を全高次構造を保つ comonad 同値へ持ち上げることができる。
命題\ref{cmo:prop:spherical-comonad-model}の有限生成対象と coKleisli 代数により、
主張した $P$ と制限 intertwiner が得られる。
\end{proof}

この系は実際の比較核の構成を、境界上の一つの自然変換の比較へ還元する。
同じ rotation を持つことだけでは
\eqref{cmo:eq:spherical-unit-matching}を満たしたことにはならない。
次の小節では非可換側の $\nu$ を中心変数の Bockstein と同定し、
Heisenberg 作用と両立する実境界核によってその像を計算する。
これにより必要な射の一致を証明し、
定理\ref{cmo:thm:fixed-zero-commutative-kernel}で指定固定値の比較を完成する。

% !TEX root = ../../mixed-sklyanin-cubics.tex
\subsection{境界の変形類と Heisenberg 作用}
\label{cmo:subsec:spherical-boundary-class}

球面制限の twist 射を、中心変数による変形の具体的な類として記述する。
この射の比較は、一般点の核を特殊化する方法とは独立の問題である。
ここでは $h=1$、$c=0$ を固定する。
$B_0=S(1,-1,-1)\otimes S(1,-1,-1)$、
$w_0=\Omega_1+100\Omega_2$、$S=B_0/(w_0)$ と置く。
$S$ とその残余圏は係数 $\alpha$ に依存しない。

\begin{lemma}[中心変数の Bockstein と twist 射]
\label{cmo:lem:spherical-bockstein}
中心正則元 $v$、$\deg v=2$ を持つ $R$ と $S=R/(v)$ に対して、
導来制限の余単位三角の射
\[
 \nu_N:N\longrightarrow N(-2)[2]
\]
は、$R/(v^2)$ 上への持上げの Bockstein である。
$S$-自由解消の微分 $d$ を $R$ 上へ持ち上げ、
$\widetilde d^{\,2}=v\theta$ と書けば、$\theta\bmod v$ がこの射を表す。
残余圏への射影と三次周期によって、この射は
命題\ref{cmo:prop:spherical-restriction-twist}の
$\nu:\operatorname{id}\to\varrho_S$ に移る。

可換側 $R^{\mathrm{cl}}=\C[x_1,\ldots,x_6,v]/(G-vL)$ では、
この射は
\begin{equation}\label{cmo:eq:commutative-twist-linear-class}
 \nu^{\mathrm{cl}}=L:
 \operatorname{id}_{\MF^{\mathbb Z}(\C[x],G)}
       \longrightarrow (1)
\end{equation}
である。ここでは三次周期 $(3)\simeq[2]$ の規約を固定している。
\end{lemma}

\begin{proof}
$R(-2)\xrightarrow{v}R$ による $S$ の解消を使い、
$S$-自由解消の各項を自由 $R$-加群へ持ち上げる。
微分の二乗の零性を $R$ 上で測ると、それは $v$ の倍数である。
三項部分の二つの短完全列を貼り合わせると、
余単位の cone に現れる二拡張の代表は $\widetilde d^{\,2}/v$ となる。
持上げを $\widetilde d+vb$ に変えると、剰余での代表の差は
$db+bd$ である。従ってその類は選択に依存せず、解消の比較射とも両立する。
これは中心元に対する cohomology operator の構成でもある
\cite[Construction~4.1.1]{CMOEisenbudPeeva}。
中心性と自由解消だけを用いるので、この記述は非可換 $R$ にも成立する。

残余圏では全圏の余単位三角へ例外射影を施す。
式\eqref{cmo:eq:spherical-twist-mutation-triangle}はこの同じ三角であるから、
その射は上記 Bockstein の射影である。
同じ次数付き自由解消と比較射を qgr 商と Orlov window に通しているので、
ここで別の自然変換へ取り替えていない。
MF 表示で内部 shift を使うと $(-2)[2]=(1)$ となる。
可換側の有限自由行列分解の微分を $d_N$ とすれば、
多項式環上で $d_N^2=G\operatorname{id}$、$R^{\mathrm{cl}}$ 上で
$d_N^2=vL\operatorname{id}$ である。
従って cohomology operator は $L\operatorname{id}$ であり、
式\eqref{cmo:eq:commutative-twist-linear-class}を得る。
\end{proof}
\begin{lemma}[六つの係数の非零な境界類]
\label{cmo:lem:boundary-parameter-injection}
$\alpha=(\alpha^{(1)},\alpha^{(2)})\in\C^6$ による中心 $v$ 変形は線形写像
\begin{equation}\label{cmo:eq:boundary-parameter-map}
 b:\C^6\longrightarrow
 H^0\operatorname{RHom}_{\operatorname{Fun}(\mathcal K_S,\mathcal K_S)}
       (\operatorname{id},\varrho_S)
\end{equation}
を定める。この写像は単射であり、実際には同型である。
固定係数での $b(\alpha)$ は上の twist 射 $\nu$ である。
\end{lemma}

\begin{proof}
$R/(v^2)$ の関係式と中心三次の $v$ 係数は $\alpha$ に線形である。
補題\ref{cmo:lem:spherical-bockstein}による変形類の構成もこの一次係数に線形なので、
写像 $b$ を得る。

単射性は剰余体の安定化上で検査できる。
$F=(S(1,-1,-1)^!)^{\widehat\otimes2}$ とし、
$\xi_{ij}$ を二次双対の六つの生成元とする。
$B_0$ の Koszul 解消は内部次数と homological 次数が等しい。
$w_0$ に対する Shamash 解消では、$n$ 番目の高次 homotopy の係数の次数は
$3n-(2n-1)=n+1>0$ である。
従ってこの解消は最小であり、その自己 Ext は、周期作用素 $t$ に関する加群として
\[
 \underline{\operatorname{Ext}}^*_S(\C,\C)
       \simeq F\otimes\C[t],\qquad
       \deg_{\mathrm{coh}}t=2,\quad \deg_{\mathrm{int}}t=3
\]
となる。この式は積全体の同型を仮定せず、最小解消と周期作用素の作用の記述である。
特に $t$ による torsion はない。

次数二の関係を $R$ へ持ち上げると、各ブロックについて
\[
 x_i^2-y_iz_i+z_iy_i=\alpha^{(i)}_1v
\]
および巡回した二式を得る。
従って Bockstein の剰余体上の類は、全体の符号の取り方を除いて
\begin{equation}\label{cmo:eq:boundary-bockstein-dual-class}
 \sum_{i,j}\alpha_j^{(i)}\xi_{ij}^{\,2}
 \in F^2\subset
 \underline{\operatorname{Ext}}^2_S(\C,\C(-2))
\end{equation}
である。中心三次の $v$ 係数は一次式なので、剰余体への augmentation 後には
この内部次数二の成分へ寄与しない。
六つの $\xi_{ij}^2$ は $F^2$ で一次独立である。
安定 Hom はこの Shamash 解消の周期的な尾部、すなわち $t$ の可逆化で計算される。
$t$-torsion がないので、非零の上の類は安定化後にも非零である。
具体的には、$t$ を十分多く掛けて $S$ の injective dimension 五を超える
解消の尾部へ移しても、この類は零にならない。
従って $b(\alpha)=0$ なら $\alpha=0$ となる。

既存の rotation を保つ分離比較で
$\mathcal K_S\simeq\Ku(V(G))$ と同定する。
滑らかな三次四様体の対角 MF の恒等 sector では、
$\operatorname{Hom}(\operatorname{id},(1))=\operatorname{Jac}(G)_1$ である。
非恒等 sector の条件はここでは満たされない
（命題\ref{cmo:prop:cubic-boundary-comonad-vanishing}の計算を参照）。
Jacobian イデアルは次数二で生成されるので、この空間は六次元である。
従って単射 $b$ は同型である。
\end{proof}

\paragraph{対称性と比較核の選択.}
境界の二ブロックには、各々位数二十七の Heisenberg 群が作用する。
その生成元は
\[
 c(x,y,z)=(y,z,x),\qquad d(x,y,z)=(x,\omega y,\omega^2z)
\]
である。普遍環の係数座標には双対作用を用いるが、
固定ファイバーのパラメータ点にはその逆変換が作用する。
例えばパラメータ点上では
\[
 c_\alpha(\alpha_1,\alpha_2,\alpha_3)
       =(\alpha_3,\alpha_1,\alpha_2),\qquad
 d_\alpha(\alpha_1,\alpha_2,\alpha_3)
       =(\alpha_1,\omega\alpha_2,\omega^2\alpha_3).
\]
中心 $x\mapsto\zeta x$ はパラメータ点を
$\alpha\mapsto\zeta^{-2}\alpha=\zeta\alpha$ へ移す。
この固定ファイバー間の pullback による規約を両側で共通に用いる。
写像\eqref{cmo:eq:boundary-parameter-map}と
$\alpha\mapsto\alpha_1x+\alpha_2y+\alpha_3z$ はともにこの作用に自然である。
境界比較を、この作用と rotation の双方に両立するように選ぶ点が必要となる。

\begin{lemma}[Heisenberg 作用と両立する楕円因子の比較]
\label{cmo:lem:heisenberg-boundary-comparison}
本例の二つの同一楕円因子について、分離比較
$E:\mathcal K_S\xrightarrow{\sim}\mathcal K_{S^{\mathrm{cl}}}$ を、
両ブロックの Heisenberg 作用と rotation に両立する実 perfect 核として選べる。
二つの中心三次の係数を等しく正規化した表示では、この選択はブロック交換にも両立する。
\end{lemma}

\begin{proof}
一つの楕円因子を $(E_0,L_0,t_p)$ とし、$G_3=E_0[3]$ と書く。
次数三の線束 $L_0$ の theta 群の有限 Heisenberg 部分が、上の線型作用を与える。
楕円 window 上の二つの rotation は
\[
 R_p=T_{\OO_{E_0}}(t_p)_*(-\otimes L_0),\qquad
 R_0=T_{\OO_{E_0}}(-\otimes L_0)
\]
である。既存の Picard 補正
\paperref{main}{mbrc:lem:factor-correction}を用い、
$[M]=p$ の実次数零線束と $F_M=-\otimes M$、
$F_MR_p\simeq R_0F_M$ を取る。

$[3]^*:\operatorname{Pic}^0(E_0)\to\operatorname{Pic}^0(E_0)$ は全射なので、
$[3]^*N\simeq M$ となる $N$ を選べる。
この表示によって $M$ に $G_3$ の線型化を与え、$F_M$ を $G_3$-同変にする。
二つの rotation の Heisenberg 中心は同じ重みで作用するため、
その間の intertwiner の一次元空間には $G_3$ の文字 $\chi$ が作用する。
ここで同型の空間が一次元であることは $HH^0(E_0)=\C$ による。

$q\in G_3$ による標的の平行移動 $(t_q)_*$ を合成する。
theta 群の lift を用いる
$(t_q)_*R_0\simeq R_0(t_q)_*$ の $G_3$ に関する文字は、
theta 群の交換子 $e_3(q,-)$ である。
ここでは pushforward $(t_q)_*$ の lift の交換子を $e_3(q,-)$ とする規約を用い、
逆の lift の規約では以下の $q$ を $-q$ に置き換える。
theta 群とその非退化な交換子については
\cite[Sections~1.3 and~1.6, Definition~1.29]{NFMCremonaFisherDescentAlgebra}を参照する。
Weil pairing $e_3$ の非退化性により、$\chi e_3(q,-)=1$ となる $q$ を選べる。
従って $F_i=(t_q)_*F_M$ と rotation の intertwiner はともに
Heisenberg 作用と両立する。
三乗周期の差は scalar なので、intertwiner にその三乗根を掛けて同期できる。
この変更は同変性を壊さない。

二つの同一因子には同じ $M,N,q$ と同じ周期補正を用いる。
外積、指定中心群に関する同変化、二重同変化を用いる
主稿の分離比較の構成は、これらの同変自然変換を保つ。
実際、相対 rotation $Q=\varrho_1\varrho_2^{-1}$ の intertwiner は
二つの因子の intertwiner の積であり、両 Heisenberg 作用に不変である。
先に同期した三乗周期によって $Q^3$ の coherence も一致する。
因子交換は $Q$ を $Q^{-1}$ へ移すので、相対中心群の同変化と
二重同変化の平均冪等元にも作用が降りる。
従って得られる境界比較は両 Heisenberg 作用、全体の rotation、
および因子交換に両立する。
中心三次の係数 $1,100$ は、第二ブロックの生成元を $100^{1/3}$ 倍して
係数 $1,1$ へ移してからこの構成を行い、元へ戻せばよい。
全ての選択は $\C$ 上で行い、有理数体上の核降下は主張しない。
\end{proof}

\begin{proposition}[twist 射の比較]
\label{cmo:prop:spherical-boundary-unit-match}
補題\ref{cmo:lem:heisenberg-boundary-comparison}の境界比較と
rotation の同型を適切に選ぶと、
\begin{equation}\label{cmo:eq:spherical-boundary-unit-match}
 E\nu=\nu^{\mathrm{cl}}E
\end{equation}
が成立する。ここでは比較対象の二つの twist をその rotation と同定している。
\end{proposition}

\begin{proof}
まず中心三次の係数を $1,1$ に正規化する。
六次元空間 $\operatorname{Hom}(\operatorname{id},\varrho)$ は、
二つの Heisenberg 群に関して、各ブロックの共通の三次元既約表現の直和である。
この二表現は互いに非同型で、各々既約である。
補題\ref{cmo:lem:boundary-parameter-injection}と前補題の同変性から、
二つのパラメータ写像の比較は各ブロック上で scalar となる。
因子交換との両立性により二つの scalar は同じ $c\in\C$ である。
パラメータ写像の単射性により $c\ne0$ である。
元の係数へ戻すと、第二ブロックの $\alpha$ は $100^{2/3}$ 倍となるので、
得られる一次式は $cL$、
$L=\alpha^{(1)}\cdot x^{(1)}+100\alpha^{(2)}\cdot x^{(2)}$ である。

従って境界比較による $\nu$ の像は $c\nu^{\mathrm{cl}}$ となる。
twist を比較する自然同型を $c^{-1}$ 倍すれば
式\eqref{cmo:eq:spherical-boundary-unit-match}を得る。
この最後の調整は、三乗周期まで標識した境界比較という追加条件を課すものではない。
本稿でここから構成するのは通常の $\C$-線型 dg Morita 同値である。
\end{proof}

\subsection{零係数の混合五様体に対する実比較核}
\label{cmo:subsec:fixed-zero-kernel}

これまでの比較を合成すると、固定値での結論を得る。
その構成では、形式的な $h$ 級数に $h=1$ を代入しない。

\begin{theorem}[指定された零係数模型の可換 Kuznetsov 同値]
\label{cmo:thm:fixed-zero-commutative-kernel}
$\alpha^{(1)}=(1,4,9)$、$\alpha^{(2)}=(16,25,49)$ とし、
$A_0$ を本稿の指定混合五様体の零係数模型とする。
命題\ref{cmo:prop:zero-eckardt-resolution}の滑らかな可換三次五様体
\[
 X_0=V(G-u^2L)\subset\mathbb P^6
\]
に対し、実際の perfect 核を持つ $\C$-線型 dg Morita 同値
\begin{equation}\label{cmo:eq:fixed-zero-commutative-kernel}
 \operatorname{Ku}_{\mathrm{coh}}(A_0)
      \xrightarrow{\ \sim\ }\operatorname{Ku}(X_0)
\end{equation}
を構成できる。

その中間段階として、第二成分の実比較
\[
 P_R:\mathcal K_R\xrightarrow{\sim}\mathcal K_{R^{\mathrm{cl}}}
       \simeq\operatorname{Ku}(Y),\qquad Y=V(G,L),
\]
および境界の分離比較 $E$ に対する dg 自然同型
\begin{equation}\label{cmo:eq:fixed-zero-restriction-square}
 E\mathcal L\xrightarrow{\ \sim\ }\mathcal L^{\mathrm{cl}}P_R
\end{equation}
を同時に構成できる。
これは通常の dg Morita 同値についての主張であり、
全ての内部 rotation・指定周期・元の中心作用の標識を同時に保存するという
追加条件は課していない。
\end{theorem}

\begin{proof}
源の全 coherent 残余圏の smooth proper 性は
系\ref{cmo:cor:fixed-zero-norm-smooth}による。
$X_0,Z,Y$ の滑らかさは命題\ref{cmo:prop:zero-eckardt-resolution}で確認済みである。
従って両側の球面制限と、その随伴の perfect 核が存在する。
命題\ref{cmo:prop:spherical-restriction-twist}により、両側の twist は境界の rotation、
cotwist の三乗は $[-1]$ である。

補題\ref{cmo:lem:heisenberg-boundary-comparison}の実境界核を選び、
命題\ref{cmo:prop:spherical-boundary-unit-match}により二つの twist 射を同定する。
系\ref{cmo:cor:spherical-unit-reconstruction}を適用して
$P_R$ と式\eqref{cmo:eq:fixed-zero-restriction-square}を得る。
この段階の核を具体的に記述すると、境界の生成対象 $G_S$ に対して
\[
 G_R^{\mathrm{adj}}=\mathcal L^RG_S,\qquad
 G_{\mathrm{cl}}^{\mathrm{adj}}
      =(\mathcal L^{\mathrm{cl}})^R E G_S
\]
を取る。これらは補題\ref{cmo:lem:periodic-spherical-generation}により生成対象である。
全 coherence を含む comonad 比較は
\[
 R\operatorname{End}(G_R^{\mathrm{adj}})
       \simeq R\operatorname{End}(G_{\mathrm{cl}}^{\mathrm{adj}})
\]
を積と高次構造を保つ準同型として与える。
その graph 双加群と二つの生成対象の Morita 核を tensor したものが $P_R$ である。
graph は体上 perfect で、両自己準同型代数は smooth proper なので、
この双加群も perfect である。準同型の逆と生成性から unit と counit は同型になる。
可換側の $\operatorname{Ku}(Y)$ への同定には
命題\ref{cmo:prop:commutative-invariant-knorrer}の実 rank 一核を使う。

最後に命題\ref{cmo:prop:root-order-comparison-criterion}を適用する。
二つの対角核 $E,P_R$ と制限の intertwiner を三角行列の双加群として貼り合わせ、
\[
 \widehat P:\operatorname{Ku}_{\mathrm{coh}}(A_0)^{\mu_2}
       \xrightarrow{\sim}\operatorname{Ku}(X_0)^{\mu_2}
\]
を得る。指標作用は Serre 関手から内在的に回復される。
その比較の二乗を単元の平方根で補正し、
明示的な平均冪等元によって同変 perfect 核を作る。
二重 crossed product の行列代数との Morita 同値を合成すれば、
式\eqref{cmo:eq:fixed-zero-commutative-kernel}の核が得られる。
これらの操作の unit と counit は、両成分の生成対象上で同型であり、
有限 cone と直和因子を通じて全残余圏で同型となる。
\end{proof}

\begin{remark}[構成の範囲]
\label{cmo:rem:fixed-zero-kernel-scope}
定理の核は、楕円曲線上の線束・平行移動核、coKleisli 代数の graph、
有限の貼合せと同変冪等元から構成した perfect 導来核である。
全行列係数を単一の有限行列として列挙したという意味ではない。
構成は $\C$ 上で線束と根を選んでおり、$\mathbb Q$ 上への核降下を結論しない。
また、この固定値の核が
定理\ref{mix:thm:h-exact-candidate}の選択済み一般点核の特殊化に一致することも
主張しない。
中心係数 $c\ne0$ の指定候補への比較、定常 norm の全六十四根を含む
全 smooth locus の比較、一般の混合表示と全 intrinsic 対象への拡張は残る。
\end{remark}

